\documentclass[12pt,a4paper]{article}
\usepackage{fontspec}
\usepackage{xltxtra}
\usepackage{geometry}
\usepackage{titlesec}
\usepackage{enumitem}
\usepackage{longtable}
\usepackage{array}
\usepackage{booktabs}
\usepackage{xcolor}
\usepackage[colorlinks=true,linkcolor=blue,urlcolor=blue]{hyperref}
\usepackage{fancyhdr}

% ฟอนต์ทางการ: TH SarabunPSK ตามมติคณะรัฐมนตรีที่กำหนดให้ใช้เป็นฟอนต์มาตรฐานราชการไทย
% (หากเครื่องที่ใช้คอมไพล์ไม่มีฟอนต์นี้ติดตั้งไว้ ให้ติดตั้งจากชุดฟอนต์ราชการ "TH SarabunPSK"
% หรือ "TH Sarabun New" ก่อนคอมไพล์)
\setmainfont[
  BoldFont={TH SarabunPSK Bold},
  ItalicFont={TH SarabunPSK Italic},
  BoldItalicFont={TH SarabunPSK Bold Italic}
]{TH SarabunPSK}

% การตัดบรรทัดภาษาไทยที่ถูกต้องสำหรับ XeLaTeX (ภาษาไทยไม่มีช่องว่างระหว่างคำ
% หากไม่กำหนด locale นี้ ข้อความจะล้นขอบกระดาษและอ่านไม่ออก)
\XeTeXlinebreaklocale "th"
\XeTeXlinebreakskip 0pt plus 1pt

\geometry{a4paper, left=2cm, right=3cm, top=2cm, bottom=2cm}

\renewcommand{\baselinestretch}{1.25}

\titleformat{\section}{\bfseries\Large}{\thesection.}{0.5em}{}
\titleformat{\subsection}{\bfseries\large}{\thesubsection}{0.5em}{}
\titleformat{\subsubsection}{\bfseries\normalsize}{}{0em}{}

\pagestyle{fancy}
\setlength{\headheight}{18pt}
\fancyhf{}
\fancyhead[L]{กฎหมายภาษีอากรในชีวิตประจำวัน}
\fancyhead[R]{ค่าลดหย่อนภาษีเงินได้บุคคลธรรมดา}
\fancyfoot[C]{\thepage}

\setlength{\parindent}{1.5em}
\setlength{\parskip}{0.4em}
\raggedright

\newcommand{\hilight}[1]{\textbf{#1}}
\newcommand{\lawref}[1]{\textit{(#1)}}

\title{\Huge กฎหมายภาษีอากรในชีวิตประจำวัน \\[0.3em] \Large ค่าลดหย่อนภาษีเงินได้บุคคลธรรมดา \\[0.2em] \large (สรุปเนื้อหาฉบับละเอียดเพื่อใช้ทบทวนและออกข้อสอบ)}
\author{}
\date{}

\begin{document}
\maketitle
\thispagestyle{fancy}

\renewcommand{\contentsname}{สารบัญ}
\tableofcontents
\newpage

\section{โครงสร้างการคำนวณภาษีเงินได้บุคคลธรรมดา}

การคำนวณภาษีเงินได้บุคคลธรรมดาตามประมวลรัษฎากรมีขั้นตอนหลักสามขั้นตอนต่อเนื่องกัน ดังนี้

\begin{enumerate}
    \item \textbf{เงินได้พึงประเมิน} ตามมาตรา 40(1)--(8) ประกอบมาตรา 39 และมาตราอื่นที่เกี่ยวข้อง ซึ่งเป็นเงินได้ทุกประเภทที่กฎหมายกำหนดให้ต้องนำมาคำนวณเพื่อเสียภาษี
    \item \textbf{หักค่าใช้จ่าย} ตามมาตรา 42 ทวิ ถึงมาตรา 46 ประกอบพระราชกฤษฎีกาฉบับที่ 11 เพื่อให้ได้เงินได้หลังหักค่าใช้จ่าย
    \item \textbf{หักค่าลดหย่อน} ตามมาตรา 47 ประกอบกฎกระทรวงฉบับที่ 126 และกฎหมายอื่นที่เกี่ยวข้อง เพื่อให้ได้ \hilight{เงินได้สุทธิ} ซึ่งจะนำไปคำนวณภาษีตามอัตราภาษีก้าวหน้าต่อไป
\end{enumerate}

มาตรา 47 แห่งประมวลรัษฎากรบัญญัติว่า

\begin{quote}
\textit{"เงินได้พึงประเมินตามมาตรา 40 เมื่อได้หักตามมาตรา 42 ทวิ ถึงมาตรา 46 แล้ว เพื่อเป็นการบรรเทาภาระภาษี ให้หักลดหย่อนได้อีกต่อไปนี้..."}
\end{quote}

บทบัญญัตินี้แสดงให้เห็นชัดเจนว่าค่าลดหย่อนเป็นมาตรการที่กฎหมายกำหนดขึ้นเพื่อ \hilight{บรรเทาภาระภาษี} ให้แก่ผู้มีเงินได้ โดยลดฐานภาษีลงก่อนนำไปคำนวณภาษีตามอัตราที่กำหนด

\subsection*{ข้อสังเกตสำคัญ: True Exemption กับ Exemption as Deduction}

ในการศึกษากฎหมายภาษีอากรเรื่องค่าลดหย่อน จำเป็นต้องแยกแยะแนวคิดสองประการที่มีลักษณะคล้ายคลึงกันแต่มีที่มาทางกฎหมายแตกต่างกัน ได้แก่

\begin{itemize}
    \item \hilight{"ยกเว้นเงินได้ที่ได้รับ" (True Exemption)} หมายถึงกรณีที่กฎหมายยกเว้นมิให้นำเงินได้ประเภทหนึ่งมารวมคำนวณเพื่อเสียภาษีตั้งแต่ต้น โดยไม่ขึ้นอยู่กับว่าผู้มีเงินได้จะใช้จ่ายเงินนั้นไปอย่างไร
    \item \hilight{"ยกเว้นเงินได้เท่าที่ได้จ่าย" (Exemption as Deduction)} หมายถึงกรณีที่กฎหมายให้สิทธิยกเว้นภาษีเงินได้เท่ากับจำนวนที่ผู้มีเงินได้ได้จ่ายไปจริงสำหรับวัตถุประสงค์ที่กฎหมายกำหนด ซึ่งมีผลในทางปฏิบัติเสมือนเป็นการหักค่าลดหย่อนอย่างหนึ่ง แม้รูปแบบทางกฎหมายจะออกมาในรูปของ "การยกเว้นภาษี" ผ่านพระราชกฤษฎีกาหรือกฎกระทรวง มิใช่การหักลดหย่อนตามมาตรา 47 โดยตรง
\end{itemize}

ตัวอย่างที่ชัดเจนของ Exemption as Deduction ได้แก่ กรณีผู้สูงอายุและคนพิการตามกฎกระทรวงฉบับที่ 126 ข้อ 2(72) และข้อ 2(81) เงินบริจาคบางประเภทที่ได้รับยกเว้น 1 เท่าหรือ 2 เท่าของจำนวนที่บริจาคตามพระราชกฤษฎีกาเฉพาะฉบับ และมาตรการกระตุ้นเศรษฐกิจประเภท "ช้อปดีมีคืน" ที่ออกตามมาตรา 42(17)

\newpage
\section{ค่าลดหย่อนส่วนตัว (มาตรา 47(1)(ก))}

หักลดหย่อนให้สำหรับผู้มีเงินได้ จำนวน \hilight{60,000 บาท}

\subsection*{หลักเกณฑ์สำคัญ}

ผู้มีเงินได้ไม่ว่าจะเป็นผู้อยู่ในประเทศไทยหรือไม่ก็ตาม ก็มีสิทธิหักค่าลดหย่อนส่วนตัวจำนวน 60,000 บาทตามมาตรา 47(1)(ก) เสมอ กล่าวคือ \hilight{ค่าลดหย่อนส่วนตัวไม่มีเงื่อนไขเรื่องถิ่นที่อยู่} ซึ่งแตกต่างจากค่าลดหย่อนคู่สมรสและบุตรที่มีเงื่อนไขจำกัดในกรณีผู้มีเงินได้มิได้เป็นผู้อยู่ในประเทศไทย (ตามมาตรา 47(3))

\subsection*{หน่วยภาษีเงินได้บุคคลธรรมดาที่มีลักษณะพิเศษ}

นอกจากบุคคลธรรมดาทั่วไป ประมวลรัษฎากรยังกำหนดหน่วยภาษีพิเศษอีก 3 กรณีที่มีสิทธิหักค่าลดหย่อนตามมาตรา 47(1)(ก) ด้วยหลักเกณฑ์เฉพาะ ได้แก่

\begin{enumerate}
    \item กรณีผู้ถึงแก่ความตายระหว่างปีภาษี --- ดูมาตรา 47(4)
    \item กรณีกองมรดกที่ยังไม่ได้แบ่ง --- ดูมาตรา 47(5)
    \item กรณีห้างหุ้นส่วนสามัญหรือคณะบุคคลที่มิใช่นิติบุคคล --- ดูมาตรา 47(6)
\end{enumerate}

(รายละเอียดของทั้งสามกรณีนี้จะอธิบายในหัวข้อที่ 12--14 ของเอกสารนี้)

\subsection*{กรณีผู้มีเงินได้เป็นผู้สูงอายุหรือคนพิการ}

นอกจากค่าลดหย่อนส่วนตัว 60,000 บาทแล้ว กฎกระทรวงฉบับที่ 126 ยังให้สิทธิประโยชน์เพิ่มเติมในรูปแบบการยกเว้นภาษีเงินได้สำหรับผู้สูงอายุและคนพิการ ดังนี้

\textbf{ข้อ 2(72)} เงินได้ที่ผู้มีเงินได้ซึ่งเป็นผู้อยู่ในประเทศไทยและมีอายุไม่ต่ำกว่า 65 ปีบริบูรณ์ในปีภาษีได้รับ ได้รับยกเว้นภาษีเฉพาะส่วนที่ไม่เกิน \hilight{190,000 บาท} ในปีภาษีนั้น สำหรับเงินได้ที่ได้รับตั้งแต่วันที่ 1 มกราคม พ.ศ. 2548 เป็นต้นไป โดยเป็นไปตามหลักเกณฑ์ วิธีการ และเงื่อนไขที่อธิบดีกรมสรรพากรประกาศกำหนด (ดูประกาศอธิบดีกรมสรรพากรเกี่ยวกับภาษีเงินได้ ฉบับที่ 150)

\textbf{ข้อ 2(81)} เงินได้ที่ผู้มีเงินได้เป็นคนพิการที่มีบัตรประจำตัวคนพิการตามกฎหมายว่าด้วยการส่งเสริมและพัฒนาคุณภาพชีวิตคนพิการ ซึ่งเป็นผู้อยู่ในประเทศไทยและมีอายุไม่เกิน 65 ปีบริบูรณ์ในปีภาษีได้รับ เฉพาะส่วนที่ไม่เกิน \hilight{190,000 บาท} สำหรับปีภาษีนั้น สำหรับเงินได้พึงประเมินที่ได้รับตั้งแต่วันที่ 1 มกราคม พ.ศ. 2553 เป็นต้นไป (แก้ไขเพิ่มเติมโดยกฎกระทรวงฉบับที่ 281 พ.ศ. 2554 ใช้บังคับวันที่ 9 พฤษภาคม 2554 เป็นต้นไป) (ดูประกาศอธิบดีกรมสรรพากรเกี่ยวกับภาษีเงินได้ ฉบับที่ 197)

\textbf{เงื่อนไขร่วมของทั้งสองกรณี}

\begin{itemize}
    \item ต้องเป็นบุคคลธรรมดา
    \item ต้องเป็นผู้อยู่ในประเทศไทย
    \item สามารถเลือกใช้สิทธิยกเว้นภาษีเงินได้จากเงินได้ที่ได้รับประเภทใดประเภทหนึ่ง หรือจะเลือกใช้สิทธิจากเงินได้หลายประเภทก็ได้ และแต่ละประเภทจะยกเว้นภาษีเงินได้จำนวนเท่าใดก็ได้ แต่เมื่อรวมจำนวนเงินได้ที่ใช้สิทธิยกเว้นภาษีเงินได้ดังกล่าวทั้งหมดแล้ว ต้องไม่เกิน \hilight{190,000 บาท} ในปีภาษีนั้น
    \item \hilight{สิทธิตามข้อ 2(72) และข้อ 2(81) ใช้ร่วมกันในปีภาษีเดียวกันไม่ได้} กล่าวคือ ผู้มีเงินได้ที่เข้าเกณฑ์ทั้งสองกรณี (เช่น อายุเกิน 65 ปี และเป็นคนพิการด้วย) ต้องเลือกใช้สิทธิเพียงข้อเดียวในปีภาษีหนึ่ง ๆ
\end{itemize}

\newpage
\section{ค่าลดหย่อนคู่สมรส (มาตรา 47(1)(ข))}

หักลดหย่อนให้สำหรับสามีหรือภริยาของผู้มีเงินได้ จำนวน \hilight{60,000 บาท}

\subsection*{เงื่อนไขพื้นฐาน}

ต้องเป็นสามีหรือภริยา \hilight{โดยชอบด้วยกฎหมาย} เท่านั้น และมีได้สูงสุดเพียง 1 คน (กรณีมีคู่สมรสหลายคนตามกฎหมายเดิมก่อนการบังคับใช้ประมวลกฎหมายแพ่งและพาณิชย์ บรรพ 5 ฉบับปัจจุบัน ก็ยังคงหักลดหย่อนได้เพียง 1 คนเท่านั้น)

\subsection*{กรณีสามีภริยาต่างฝ่ายต่างมีเงินได้ (มาตรา 47(2) ประกอบมาตรา 57 ฉ)}

ในกรณีสามีและภริยาต่างฝ่ายต่างมีเงินได้ การหักค่าลดหย่อนตามมาตรา 47(1)(ก) [ส่วนตัว] และมาตรา 47(1)(ข) [คู่สมรส] ให้หักลดหย่อนรวมกันได้ \hilight{120,000 บาท} เสมอ ไม่ว่าทั้งสองฝ่ายจะเลือกวิธียื่นแบบแสดงรายการในรูปแบบใดก็ตาม

\subsection*{วิธีการยื่นแบบแสดงรายการของสามีภริยาตามมาตรา 57 ฉ}

มาตรา 57 ฉ แห่งประมวลรัษฎากรบัญญัติว่า

\begin{quote}
{\itshape "ในการเก็บภาษีเงินได้จากสามีและภริยานั้น ให้สามีและภริยาต่างฝ่ายต่างมีหน้าที่ยื่นรายการเกี่ยวกับเงินได้พึงประเมินที่ตนได้รับในระหว่างปีภาษีที่ล่วงมาแล้วตามมาตรา 56

ในกรณีที่เงินได้พึงประเมินไม่อาจแยกได้อย่างชัดแจ้งว่าเป็นของสามีหรือภริยาแต่ละฝ่ายจำนวนเท่าใด ให้ถือเป็นเงินได้พึงประเมินของสามีและภริยาฝ่ายละกึ่งหนึ่ง เว้นแต่เงินได้พึงประเมินตามมาตรา 40(8) สามีและภริยาจะแบ่งเงินได้พึงประเมินเป็นของแต่ละฝ่ายตามส่วนที่ตกลงกันก็ได้ แต่รวมกันต้องไม่น้อยกว่าเงินได้พึงประเมินที่ได้รับ ถ้าตกลงกันไม่ได้ ให้ถือเป็นเงินได้พึงประเมินของสามีและภริยาฝ่ายละกึ่งหนึ่ง

สามีและภริยาจะตกลงยื่นรายการและเสียภาษีรวมกัน โดยให้ถือเอาเงินได้พึงประเมินของตนเป็นเงินได้ของสามีหรือภริยาอีกฝ่ายหนึ่งก็ได้ หรือจะแยกยื่นรายการและเสียภาษีเฉพาะส่วนที่เป็นเงินได้พึงประเมินตามมาตรา 40(1) โดยมิให้ถือเอาเป็นเงินได้ของอีกฝ่ายหนึ่งก็ได้ แต่ถ้ามีภาษีค้างชำระ สามีและภริยาต้องร่วมรับผิดในการเสียภาษีที่ค้างชำระนั้น

เมื่อได้เลือกยื่นรายการตามวรรคสองและวรรคสามในปีภาษีใดแล้ว ให้ถือว่าเป็นวิธีการยื่นรายการสำหรับปีภาษีนั้นตลอดไป เว้นแต่อธิบดีจะอนุมัติให้เปลี่ยนแปลงวิธีการเลือกยื่นรายการดังกล่าว"}
\end{quote}

จากบทบัญญัติดังกล่าว สามีภริยาสามารถเลือกวิธียื่นแบบแสดงรายการได้ 3 รูปแบบ ดังนี้

\textbf{แบบที่ 1 (แยกยื่น)} สามีและภริยาต่างฝ่ายต่างมีหน้าที่ยื่นรายการเกี่ยวกับเงินได้พึงประเมินที่ตนได้รับ ในกรณีที่เงินได้พึงประเมินไม่อาจแยกได้อย่างชัดแจ้งว่าเป็นของฝ่ายใดเท่าใด ให้ถือเป็นเงินได้พึงประเมินของสามีและภริยาฝ่ายละกึ่งหนึ่ง เว้นแต่เงินได้พึงประเมินตามมาตรา 40(8) สามีภริยาจะแบ่งเงินได้พึงประเมินเป็นของแต่ละฝ่ายตามส่วนที่ตกลงกันก็ได้ (แต่รวมกันต้องไม่น้อยกว่าเงินได้พึงประเมินที่ได้รับ) ถ้าตกลงกันไม่ได้ให้ถือเป็นเงินได้พึงประเมินของสามีและภริยาฝ่ายละกึ่งหนึ่ง

\textbf{แบบที่ 2 (รวมยื่น)} สามีและภริยาจะตกลงยื่นรายการและเสียภาษีรวมกัน โดยให้ถือเอาเงินได้พึงประเมินของตนเป็นเงินได้ของสามีหรือภริยาอีกฝ่ายหนึ่งก็ได้ ซึ่งมีทางเลือก 2 ทาง คือ
\begin{enumerate}
    \item ภริยารวมยื่นกับสามี (ถือเอาเงินได้ของภริยาเป็นเงินได้ของสามี)
    \item สามีรวมยื่นกับภริยา (ถือเอาเงินได้ของสามีเป็นเงินได้ของภริยา)
\end{enumerate}

\textbf{แบบที่ 3 (แยกยื่นเฉพาะเงินได้ตามมาตรา 40(1) ที่เหลือรวมยื่น)} สามีหรือภริยาจะแยกยื่นรายการและเสียภาษีเฉพาะส่วนที่เป็นเงินได้พึงประเมินตามมาตรา 40(1) (เงินเดือนค่าจ้าง) โดยมิให้ถือเอาเป็นเงินได้ของอีกฝ่ายหนึ่งก็ได้ ส่วนเงินได้ประเภทอื่นที่เหลือนำไปรวมยื่นกับอีกฝ่าย ซึ่งมีทางเลือก 2 ทาง คือ
\begin{enumerate}
    \item สามีแยกยื่นเงินได้ประเภทที่ 1 ส่วนเงินได้อื่นที่เหลือรวมยื่นกับภริยา
    \item ภริยาแยกยื่นเงินได้ประเภทที่ 1 ส่วนเงินได้อื่นที่เหลือรวมยื่นกับสามี
\end{enumerate}

\textbf{ข้อสำคัญที่ต้องพึงระวัง} หากมีภาษีค้างชำระ สามีและภริยาต้องร่วมรับผิดในการเสียภาษีที่ค้างชำระนั้นเสมอ ไม่ว่าจะเลือกวิธียื่นแบบใด และเมื่อได้เลือกยื่นรายการตามวิธีใดแบบหนึ่งในปีภาษีใดแล้ว ให้ถือว่าเป็นวิธีการยื่นรายการสำหรับปีภาษีนั้นตลอดไป เว้นแต่อธิบดีจะอนุมัติให้เปลี่ยนแปลงวิธีการเลือกยื่นรายการดังกล่าว

\subsection*{ตัวอย่างผลของมาตรา 47(2): การหักค่าลดหย่อนส่วนตัวและคู่สมรส}

\textbf{กรณีที่ 1: สามีเป็นผู้มีเงินได้ ภริยาไม่มีเงินได้}

หักค่าลดหย่อน: ส่วนตัว 60,000 บาท (มาตรา 47(1)(ก)) และคู่สมรส 60,000 บาท (มาตรา 47(1)(ข)) รวม 120,000 บาท

\textbf{กรณีที่ 2: สามีเป็นผู้มีเงินได้ ภริยาเป็นผู้มีเงินได้ (แบบที่ 2 รวมยื่น)}

ไม่ว่าจะเป็น (1) ภริยารวมยื่นกับสามี หรือ (2) สามีรวมยื่นกับภริยา ผลลัพธ์เหมือนกันคือ ฝ่ายที่เป็นผู้ยื่นหลักหักค่าลดหย่อนได้ ส่วนตัว 60,000 บาท และคู่สมรส 60,000 บาท รวม 120,000 บาท

\textbf{กรณีที่ 3: สามีเป็นผู้มีเงินได้ ภริยาเป็นผู้มีเงินได้ (แบบที่ 1 แยกยื่น)}

ทั้งสามีและภริยาต่างฝ่ายต่างหักค่าลดหย่อนส่วนตัว 60,000 บาท และคู่สมรส 60,000 บาท ฝ่ายละ 120,000 บาท (รวมทั้งคู่ 240,000 บาท)

\textbf{กรณีที่ 4: สามีเป็นผู้มีเงินได้ ภริยาเป็นผู้มีเงินได้ (แบบที่ 3 แยกยื่นเฉพาะ ม.40(1) ที่เหลือรวมยื่น)}

ไม่ว่าจะเป็น (1) สามีแยกยื่นเงินได้ประเภทที่ 1 เงินได้อื่นรวมยื่นกับภริยา หรือ (2) ภริยาแยกยื่นเงินได้ประเภทที่ 1 เงินได้อื่นรวมยื่นกับสามี ทั้งสองฝ่ายต่างหักค่าลดหย่อนส่วนตัว 60,000 บาท และคู่สมรส 60,000 บาท ฝ่ายละ 120,000 บาท เช่นเดียวกับกรณีที่ 3

\hilight{ข้อสังเกตสำคัญสำหรับการทำข้อสอบ:} ไม่ว่าสามีภริยาจะเลือกวิธียื่นแบบใดในมาตรา 57 ฉ (แยกยื่น รวมยื่น หรือแยกยื่นเฉพาะ ม.40(1)) หากทั้งสองฝ่ายต่างมีเงินได้ ทั้งสองฝ่ายจะหักค่าลดหย่อนส่วนตัวและคู่สมรสได้ฝ่ายละ 60,000 + 60,000 = 120,000 บาทเสมอ (ผลลัพธ์เท่ากันทุกวิธี) ความแตกต่างระหว่างวิธียื่นแบบจะส่งผลต่อการคำนวณภาษีในส่วนของอัตราภาษีก้าวหน้าและฐานเงินได้ที่นำมาคำนวณ มิใช่ส่งผลต่อจำนวนค่าลดหย่อนที่หักได้

\newpage
\section{ค่าลดหย่อนบุตร (มาตรา 47(1)(ค))}

\begin{itemize}
    \item บุตรชอบด้วยกฎหมาย: คนละ \hilight{30,000 บาท}
    \item บุตรบุญธรรม: คนละ \hilight{30,000 บาท}
\end{itemize}

\textbf{กรณีพิเศษ:} บุตรชอบด้วยกฎหมายตั้งแต่คนที่สองเป็นต้นไป ที่เกิดในหรือหลังปี พ.ศ. 2561 ให้หักลดหย่อนได้เพิ่มอีกคนละ \hilight{30,000 บาท} (รวมเป็น 60,000 บาทต่อคนสำหรับบุตรกลุ่มนี้) บทบัญญัตินี้แก้ไขเพิ่มเติมโดยพระราชบัญญัติแก้ไขเพิ่มเติมประมวลรัษฎากร (ฉบับที่ 46) พ.ศ. 2561 ใช้บังคับสำหรับเงินได้พึงประเมินประจำปีภาษี พ.ศ. 2561 ที่จะต้องยื่นรายการในปี พ.ศ. 2562 เป็นต้นไป

มาตรา 47(1)(ค) บัญญัติว่า

\begin{quote}
\textit{"(1) บุตรชอบด้วยกฎหมายของผู้มีเงินได้ หรือบุตรชอบด้วยกฎหมายของสามีหรือภริยาของผู้มีเงินได้ คนละ 30,000 บาท และสำหรับบุตรชอบด้วยกฎหมายตั้งแต่คนที่สองเป็นต้นไป ที่เกิดในหรือหลังปี พ.ศ. 2561 ให้หักลดหย่อนได้เพิ่มอีกคนละ 30,000 บาท โดยในการนับลำดับบุตรให้นับลำดับของบุตรทุกคนไม่ว่าจะมีชีวิตอยู่หรือไม่ก็ตาม"}
\end{quote}

\begin{quote}
\textit{"(2) บุตรบุญธรรมของผู้มีเงินได้คนละ 30,000 บาท แต่รวมกันต้องไม่เกินสามคน"}
\end{quote}

\subsection*{หลักการนับลำดับและจำนวนบุตร}

\textbf{การนับลำดับบุตร} ให้นับลำดับบุตรชอบด้วยกฎหมายทุกคน (ทั้งบุตรชอบด้วยกฎหมายของผู้มีเงินได้หรือของคู่สมรสของผู้มีเงินได้) ไม่ว่าจะมีชีวิตอยู่หรือไม่ก็ตาม

\textbf{กรณีมีบุตรทั้งบุตรชอบด้วยกฎหมายและบุตรบุญธรรม} การหักลดหย่อนสำหรับบุตรให้นำบุตรตาม (1) [บุตรชอบด้วยกฎหมาย] ทั้งหมดมาหักก่อน แล้วจึงนำบุตรตาม (2) [บุตรบุญธรรม] มาหัก เว้นแต่ในกรณีที่ผู้มีเงินได้มีบุตรตาม (1) ที่มีชีวิตอยู่ รวมเป็นจำนวนตั้งแต่สามคนขึ้นไป จะนำบุตรตาม (2) มาหักไม่ได้ แต่ถ้าบุตรตาม (1) มีจำนวนไม่ถึงสามคน ให้นำบุตรตาม (2) มาหักได้ โดยเมื่อรวมกับบุตรตาม (1) แล้วต้องไม่เกินสามคน

\textbf{การนับจำนวนบุตร} ให้นับเฉพาะบุตรที่มีชีวิตอยู่ตามลำดับอายุสูงสุดของบุตร โดยให้นับรวมทั้งบุตรที่ไม่อยู่ในเกณฑ์ได้รับการหักลดหย่อนด้วยการหักลดหย่อนสำหรับบุตร ให้หักได้เฉพาะบุตรซึ่งมีอายุไม่เกินยี่สิบห้าปีและยังศึกษาอยู่ในมหาวิทยาลัยหรือชั้นอุดมศึกษา หรือซึ่งเป็นผู้เยาว์ หรือศาลสั่งให้เป็นคนไร้ความสามารถหรือเสมือนไร้ความสามารถอันอยู่ในความอุปการะเลี้ยงดู แต่มิให้หักลดหย่อนสำหรับบุตรดังกล่าวที่มีเงินได้พึงประเมินในปีภาษีที่ล่วงมาแล้วตั้งแต่ 30,000 บาทขึ้นไป โดยเงินได้พึงประเมินนั้นไม่เข้าลักษณะตามมาตรา 42

การหักลดหย่อนสำหรับบุตรดังกล่าวให้หักได้ตลอดปีภาษีไม่ว่ากรณีที่จะหักได้นั้นจะมีอยู่ตลอดปีภาษีหรือไม่ และในกรณีบุตรบุญธรรมนั้นให้หักลดหย่อนในฐานะบุตรบุญธรรมได้แต่ฐานะเดียว (พระราชบัญญัติแก้ไขเพิ่มเติมประมวลรัษฎากร ฉบับที่ 44 พ.ศ. 2560)

\subsection*{สรุปเงื่อนไขสำคัญ 6 ประการ กรณีค่าลดหย่อนบุตร}

\begin{enumerate}
    \item \textbf{ลำดับบุตรและปีที่เกิด:} บุตรชอบด้วยกฎหมายลำดับที่ 2 เป็นต้นไปที่เกิดในหรือหลังปี พ.ศ. 2561 สามารถหักลดหย่อนได้เพิ่มอีกคนละ 30,000 บาท (30,000 + 30,000 = 60,000 บาท) การนับลำดับบุตรให้นับลำดับบุตรชอบด้วยกฎหมายทุกคน (ทั้งของผู้มีเงินได้หรือของคู่สมรสของผู้มีเงินได้) ไม่ว่าจะมีชีวิตอยู่หรือไม่ก็ตาม
    \item \textbf{จำนวนบุตร:} บุตรชอบด้วยกฎหมายสามารถหักค่าลดหย่อนได้ไม่จำกัดจำนวน ส่วนบุตรบุญธรรมหักลดหย่อนได้จำนวนไม่เกิน 3 คน ถ้ามีทั้งบุตรชอบด้วยกฎหมายและบุตรบุญธรรม การพิจารณาจำนวนบุตรให้หักบุตรชอบด้วยกฎหมายก่อน หากยังไม่ครบ 3 คนให้หักบุตรบุญธรรมได้ แต่รวมแล้วต้องไม่เกิน 3 คน การนับจำนวนบุตรให้นับเฉพาะบุตรที่ยังมีชีวิตตามลำดับอายุสูงสุดของบุตร และนับรวมทั้งบุตรที่ไม่อยู่ในเกณฑ์ได้รับการหักลดหย่อนด้วย (เช่น ไม่ผ่านเกณฑ์เงื่อนไขเรื่องอายุ ความสามารถ การศึกษา หรือเกณฑ์เงินได้พึงประเมินที่บุตรได้รับ เป็นต้น)
    \item \textbf{อายุ ความสามารถ และการศึกษาของบุตร} (ทั้งบุตรชอบด้วยกฎหมายและบุตรบุญธรรม) บุตรต้องเป็นผู้เยาว์ (อายุไม่ถึง 20 ปีบริบูรณ์) หากบุตรไม่ใช่ผู้เยาว์ บุตรต้องเป็น (ก) บุคคลที่มีอายุไม่เกิน 25 ปีและกำลังศึกษาอยู่ในระดับอนุปริญญาหรือปริญญาตรีขึ้นไป หรือ (ข) บุคคลที่ถูกศาลสั่งให้เป็นคนไร้ความสามารถหรือคนเสมือนไร้ความสามารถ
    \item บุตร (ทั้งบุตรชอบด้วยกฎหมายและบุตรบุญธรรม) ต้องอยู่ในความอุปการะเลี้ยงดูของผู้มีเงินได้
    \item เงินได้พึงประเมินซึ่งไม่ได้รับการยกเว้นภาษีที่บุตรได้รับในแต่ละปีภาษีต้อง \hilight{ต่ำกว่า 30,000 บาท}
    \item การหักค่าลดหย่อนบุตรในฐานะบุตรบุญธรรม หักได้ในฐานะบุตรบุญธรรมฐานะเดียวเท่านั้น (ไม่สามารถเปลี่ยนไปหักในฐานะอื่นได้)
\end{enumerate}

\subsection*{ข้อสังเกตสำคัญที่สุด: ไม่มีเงื่อนไขจำกัดแบบค่าลดหย่อนคู่สมรส}

การหักค่าลดหย่อนบุตร \hilight{ไม่มีเงื่อนไขจำกัด} แบบกรณีค่าลดหย่อนคู่สมรสตามมาตรา 47(2) ดังนั้นกรณีที่สามีภริยาต่างมีเงินได้ ต่างฝ่ายต่างสามารถหักค่าลดหย่อนบุตรได้ \hilight{ฝ่ายละเต็มจำนวน} สำหรับบุตรที่อยู่ในเกณฑ์ที่หักลดหย่อนได้ตามกฎหมาย ไม่ว่าสามีภริยาจะแยกยื่นแบบฯ หรือรวมยื่นแบบฯ

\subsection*{ตัวอย่างประกอบ (ใช้บุตรชอบด้วยกฎหมาย 4 คน เกิดก่อนปี 2561 และอยู่ในเกณฑ์ที่หักได้ทั้งหมด)}

\textbf{กรณีสามีเป็นผู้มีเงินได้ ภริยาไม่มีเงินได้:} หักค่าลดหย่อนส่วนตัว 60,000 บาท คู่สมรส 60,000 บาท และบุตร 4 คน 120,000 บาท (30,000 $\times$ 4)

\textbf{กรณีสามีภริยาต่างมีเงินได้ (ไม่ว่าจะแยกยื่น รวมยื่น หรือแยกยื่นเฉพาะ ม.40(1)):} ทั้งสามีและภริยาต่างฝ่ายต่างหักค่าลดหย่อนบุตร 4 คน 120,000 บาทเต็มจำนวนเช่นเดียวกัน (นอกเหนือจากค่าลดหย่อนส่วนตัวและคู่สมรสที่หักได้ฝ่ายละ 60,000 + 60,000 บาทตามที่กล่าวมาแล้ว)

\subsection*{เงื่อนไขกรณีผู้มีเงินได้มิได้เป็นผู้อยู่ในประเทศไทย (มาตรา 47(3))}

มาตรา 47(3) บัญญัติว่า

\begin{quote}
\textit{"(3) ในกรณีผู้มีเงินได้มิได้เป็นผู้อยู่ในประเทศไทย การหักลดหย่อนตาม (1)(ข) และ (ค) ให้หักได้เฉพาะสามีหรือภริยาและบุตรที่อยู่ในประเทศไทย"} (พระราชบัญญัติแก้ไขเพิ่มเติมประมวลรัษฎากร ฉบับที่ 44 พ.ศ. 2560)
\end{quote}

กล่าวคือ กรณีผู้มีเงินได้มิได้เป็นผู้อยู่ในประเทศไทย การหักลดหย่อนคู่สมรสและบุตรตามมาตรา 47(1)(ข) และ (ค) ให้หักได้ \hilight{เฉพาะกรณีที่คู่สมรสหรือบุตรเป็นผู้อยู่ในประเทศไทยเท่านั้น} ในทางตรงข้าม หากผู้มีเงินได้เป็นผู้อยู่ในประเทศไทย จะหักค่าลดหย่อนคู่สมรสและบุตรได้ทั้งกรณีที่คู่สมรสและบุตรเป็นผู้อยู่ในประเทศไทย และกรณีที่คู่สมรสและบุตรไม่ใช่ผู้อยู่ในประเทศไทย (ไม่มีเงื่อนไขจำกัด)

\textbf{ตัวอย่างที่ 1:} สามีเป็นผู้มีเงินได้และ \hilight{ไม่ใช่}ผู้อยู่ในประเทศไทย ภริยาไม่มีเงินได้แต่เป็นผู้อยู่ในประเทศไทย บุตรไม่มีเงินได้และไม่ใช่ผู้อยู่ในประเทศไทย --- กรณีนี้เข้ามาตรา 47(3) สามีหักค่าลดหย่อนคู่สมรสได้ (เพราะภริยาอยู่ในประเทศไทย) แต่หักค่าลดหย่อนบุตรไม่ได้ (เพราะบุตรไม่อยู่ในประเทศไทย) ผลคือหักได้เพียงส่วนตัว 60,000 บาท และคู่สมรส 60,000 บาท

\textbf{ตัวอย่างที่ 2:} สามีเป็นผู้มีเงินได้และ \hilight{เป็น}ผู้อยู่ในประเทศไทย ภริยาไม่มีเงินได้แต่เป็นผู้อยู่ในประเทศไทย บุตรไม่มีเงินได้และไม่ใช่ผู้อยู่ในประเทศไทย --- กรณีนี้ไม่เข้ามาตรา 47(3) สามีหักค่าลดหย่อนได้ทั้งคู่สมรสและบุตรเต็มจำนวน (ส่วนตัว 60,000 + คู่สมรส 60,000 + บุตร 30,000 บาท) แม้บุตรจะไม่อยู่ในประเทศไทยก็ตาม เพราะเงื่อนไขจำกัดตามมาตรา 47(3) ใช้เฉพาะกรณีผู้มีเงินได้ไม่อยู่ในประเทศไทยเท่านั้น

\newpage
\section{ค่าลดหย่อนค่าฝากครรภ์และค่าคลอดบุตร}

ค่าลดหย่อนนี้กำหนดไว้ในข้อ 2(99) แห่งกฎกระทรวงฉบับที่ 126 ซึ่งบัญญัติว่า

\begin{quote}
\textit{"เงินได้เท่าที่ผู้มีเงินได้หรือคู่สมรสได้จ่ายเป็นค่าฝากครรภ์และค่าคลอดบุตรของตนตามจำนวนที่จ่ายจริง สำหรับการตั้งครรภ์แต่ละคราวแต่ไม่เกินหกหมื่นบาท หากการจ่ายค่าฝากครรภ์และค่าคลอดบุตรสำหรับการตั้งครรภ์แต่ละคราวมิได้จ่ายในปีภาษีเดียวกัน ให้ได้รับยกเว้นภาษีตามจำนวนที่จ่ายจริงในแต่ละปีภาษี แต่เมื่อรวมกันแล้วต้องไม่เกินหกหมื่นบาท ทั้งนี้สำหรับเงินได้พึงประเมินที่ได้รับตั้งแต่วันที่ 1 มกราคม พ.ศ. 2561 เป็นต้นไป และให้เป็นไปตามหลักเกณฑ์ วิธีการ และเงื่อนไขที่อธิบดีประกาศกำหนด"} (ดูประกาศอธิบดีกรมสรรพากรเกี่ยวกับภาษีเงินได้ ฉบับที่ 331)
\end{quote}

\textbf{ความหมายของค่าฝากครรภ์และค่าคลอดบุตร} หมายถึง ค่าใช้จ่ายที่เกิดขึ้นจากการรักษาพยาบาลอันเนื่องมาจากการตั้งครรภ์และคลอดบุตร ไม่ว่าจะเป็นค่าตรวจและรับฝากครรภ์ ค่าบำบัดทางการแพทย์ ค่ายาและค่าเวชภัณฑ์ ค่าทำคลอด และค่ากินอยู่ในสถานพยาบาล ทั้งนี้ไม่ว่าทารกที่คลอดจะมีชีวิตรอดหรือไม่

\subsection*{สรุปเงื่อนไขต่าง ๆ กรณีค่าลดหย่อนเพื่อการฝากครรภ์และคลอดบุตร}

\begin{itemize}
    \item ลดหย่อนได้ไม่ว่าเป็นการจ่ายให้แก่สถานพยาบาลของทางราชการหรือสถานพยาบาลของเอกชน
    \item กรณีสามีหรือภริยามีเงินได้ฝ่ายเดียว ให้ยกเว้นภาษีเงินได้ให้แก่สามีหรือภริยาซึ่งเป็นผู้มีเงินได้ตามจำนวนที่จ่ายจริง แต่ไม่เกิน 60,000 บาท
    \item กรณีสามีภริยาต่างฝ่ายต่างมีเงินได้:
    \begin{itemize}
        \item ถ้ายื่นรายการรวมกัน ให้สามีหรือภริยาที่เป็นผู้ยื่นภาษีหลัก ได้รับสิทธิลดหย่อนเต็มจำนวนตามที่จ่ายจริง แต่ไม่เกิน 60,000 บาท
        \item ถ้าแยกยื่นรายการ ให้ \hilight{เฉพาะภริยา} ได้รับสิทธิลดหย่อนเต็มจำนวนตามที่จ่ายจริง แต่ไม่เกิน 60,000 บาท
    \end{itemize}
    \item สำหรับการตั้งครรภ์ในคราวเดียวกันแต่คนละปีภาษี ให้ได้รับลดหย่อนตามจำนวนที่จ่ายจริงในแต่ละปีภาษี แต่เมื่อรวมกันแล้วต้องไม่เกิน 60,000 บาท
    \item สำหรับการตั้งครรภ์หลายคราวในปีภาษีเดียวกัน ให้ได้รับลดหย่อนตามจำนวนที่จ่ายจริงสำหรับการตั้งครรภ์แต่ละคราว คราวละไม่เกิน 60,000 บาท (แยกแต่ละคราว มิใช่รวมทุกคราวแล้วจำกัดที่ 60,000 บาท)
    \item ต้องแสดงใบเสร็จรับเงินและใบรับรองแพทย์ที่ได้ตรวจและแสดงความเห็นว่าผู้มีเงินได้หรือคู่สมรสมีภาวะตั้งครรภ์
    \item ค่าลดหย่อนเพื่อการฝากครรภ์และคลอดบุตรตามข้อ 2(99) นี้ เมื่อนำไปรวมกับสิทธิการเบิกค่าฝากครรภ์และค่าคลอดบุตรซึ่งเป็นสิทธิสวัสดิการจากภาครัฐหรือภาคเอกชนที่ผู้มีเงินได้และคู่สมรสได้รับสำหรับการตั้งครรภ์แต่ละคราวแล้ว ต้องไม่เกิน 60,000 บาท
\end{itemize}

\newpage
\section{ค่าลดหย่อนเบี้ยประกันชีวิตของผู้มีเงินได้และคู่สมรส}

ค่าลดหย่อนประเภทนี้มีฐานทางกฎหมายซ้อนกันสองชั้น คือมาตรา 47(1)(ง) แห่งประมวลรัษฎากร และข้อ 2(61) แห่งกฎกระทรวงฉบับที่ 126

\subsection*{ฐานทางกฎหมายที่ 1: มาตรา 47(1)(ง) แห่งประมวลรัษฎากร}

\begin{quote}
{\itshape "(ง) เบี้ยประกันภัยที่ผู้มีเงินได้จ่ายไปในปีภาษีสำหรับการประกันชีวิตของผู้มีเงินได้ตามจำนวนที่จ่ายจริงแต่ไม่เกิน 10,000 บาท ทั้งนี้เฉพาะในกรณีที่กรมธรรม์ประกันชีวิตมีกำหนดเวลาตั้งแต่สิบปีขึ้นไป และการประกันชีวิตนั้นได้เอาประกันไว้กับผู้รับประกันภัยที่ประกอบกิจการประกันชีวิตในราชอาณาจักร

ในกรณีสามีหรือภริยาของผู้มีเงินได้มีการประกันชีวิตและความเป็นสามีภริยาได้มีอยู่ตลอดปีภาษี ให้หักลดหย่อนได้ด้วยสำหรับเบี้ยประกันที่จ่ายสำหรับการประกันชีวิตของสามีหรือภริยานั้น ตามเกณฑ์ในวรรคหนึ่ง"} (พระราชกำหนดแก้ไขเพิ่มเติม ฉบับที่ 4 พ.ศ. 2516 ใช้บังคับปีภาษี 2517 เป็นต้นไป)
\end{quote}

\subsection*{ฐานทางกฎหมายที่ 2: ข้อ 2(61) วรรคหนึ่ง แห่งกฎกระทรวงฉบับที่ 126}

\begin{quote}
\textit{"(61) เงินได้เท่าที่ผู้มีเงินได้จ่ายเป็นเบี้ยประกันภัยในปีภาษี สำหรับการประกันชีวิตของผู้มีเงินได้ตามจำนวนที่จ่ายจริง เฉพาะส่วนที่เกินหนึ่งหมื่นบาทแต่ไม่เกินเก้าหมื่นบาท โดยกรมธรรม์ประกันชีวิตต้องมีกำหนดเวลาตั้งแต่สิบปีขึ้นไป และการประกันชีวิตนั้นได้เอาประกันไว้กับผู้รับประกันภัยที่ประกอบกิจการประกันชีวิตในราชอาณาจักร ทั้งนี้ สำหรับเบี้ยประกันภัยที่ได้จ่ายตั้งแต่วันที่ 1 มกราคม พ.ศ. 2551 เป็นต้นไป"}
\end{quote}

กล่าวคือเมื่อรวมฐานทั้งสองแล้ว ผู้มีเงินได้สามารถหักค่าเบี้ยประกันชีวิตของตนเองได้ตามจำนวนที่จ่ายจริง \hilight{สูงสุดไม่เกิน 100,000 บาท} (10,000 บาทตามมาตรา 47(1)(ง) บวกส่วนเพิ่ม 90,000 บาทตามข้อ 2(61) วรรคหนึ่ง)

\subsection*{ฐานทางกฎหมายที่ 3: ข้อ 2(61) วรรคสอง --- กรณีประกันชีวิตแบบบำนาญ}

\begin{quote}
\textit{"หากเบี้ยประกันภัยที่จ่ายตามวรรคหนึ่ง เป็นเบี้ยประกันภัยสำหรับการประกันชีวิตแบบบำนาญที่จ่ายตั้งแต่วันที่ 1 มกราคม พ.ศ. 2553 เป็นต้นไป ให้เงินได้ได้รับยกเว้นไม่ต้องนำมารวมคำนวณเพื่อเสียภาษีเงินได้เพิ่มขึ้นอีกในอัตราร้อยละสิบห้าของเงินได้พึงประเมินแต่ไม่เกินสองแสนบาท ทั้งนี้ เมื่อรวมคำนวณกับเงินได้ที่ได้รับยกเว้นไม่ต้องรวมคำนวณเพื่อเสียภาษีเงินได้สำหรับกรณีที่ผู้มีเงินได้จ่ายเป็นเงินสะสมเข้ากองทุนสำรองเลี้ยงชีพตามกฎหมายว่าด้วยกองทุนสำรองเลี้ยงชีพตาม (35) หรือเงินสะสมเข้ากองทุนบำเหน็จบำนาญข้าราชการตามกฎหมายว่าด้วยกองทุนบำเหน็จบำนาญข้าราชการตาม (43) หรือเงินสะสมเข้ากองทุนสงเคราะห์ตามกฎหมายว่าด้วยโรงเรียนเอกชนตาม (54) แล้วแต่กรณี หรือเงินค่าซื้อหน่วยลงทุนในกองทุนรวมเพื่อการเลี้ยงชีพตามกฎหมายว่าด้วยหลักทรัพย์และตลาดหลักทรัพย์ตาม (55) แล้ว ต้องไม่เกินห้าแสนบาท ในปีภาษีเดียวกัน การได้รับยกเว้นตามวรรคสอง ให้เป็นไปตามหลักเกณฑ์และวิธีการที่อธิบดีกำหนด"} (แก้ไขเพิ่มเติมโดยกฎกระทรวง ฉบับที่ 279 พ.ศ. 2554 ใช้บังคับวันที่ 23 กุมภาพันธ์ 2554 เป็นต้นไป)
\end{quote}

กล่าวคือ เบี้ยประกันชีวิตแบบบำนาญได้รับยกเว้นเพิ่มเติมในอัตราร้อยละ 15 ของเงินได้พึงประเมิน แต่ไม่เกิน \hilight{200,000 บาท} ทั้งนี้เมื่อรวมกับเงินได้ที่ได้รับยกเว้นจากเงินสะสมกองทุนสำรองเลี้ยงชีพ กบข. กองทุนสงเคราะห์โรงเรียนเอกชน หรือเงินค่าซื้อหน่วยลงทุน RMF แล้ว ต้องไม่เกิน \hilight{500,000 บาท} ในปีภาษีเดียวกัน

\subsection*{สรุปเงื่อนไขต่าง ๆ กรณีค่าลดหย่อนเบี้ยประกันชีวิตระยะยาว}

\begin{itemize}
    \item \textbf{อายุกรมธรรม์:} 10 ปีขึ้นไป (กรณีประกันชีวิตแบบบำนาญ ต้องมีการกำหนดการจ่ายผลประโยชน์เงินบำนาญเป็นรายงวดอย่างสม่ำเสมอ ผู้รับประกันต้องประกอบกิจการในราชอาณาจักร โดยช่วงอายุของการจ่ายผลประโยชน์เงินบำนาญเมื่อผู้มีเงินได้มีอายุตั้งแต่ 55 ปีขึ้นไป ถึงอายุ 85 ปีหรือกว่านั้น และผู้มีเงินได้ต้องจ่ายเบี้ยประกันภัยครบถ้วนแล้วก่อนได้รับผลประโยชน์เงินบำนาญ)
    \item \textbf{กรณีมีความคุ้มครองอื่นเพิ่มเติม:} ค่าเบี้ยประกันภัยที่จ่ายสำหรับความคุ้มครองอื่นเพิ่มเติมดังกล่าว ไม่สามารถยกเว้นภาษีสำหรับเบี้ยประกันภัยดังกล่าวได้ (ได้เฉพาะส่วนของเบี้ยประกันชีวิตเท่านั้น เว้นแต่เบี้ยประกันในส่วนของความคุ้มครองอื่นมีกฎหมายกำหนดให้สิทธิประโยชน์เพิ่มเติมเป็นพิเศษ เช่น กรณีเบี้ยประกันสุขภาพของผู้มีเงินได้)
    \item \textbf{กรมธรรม์ประกันชีวิตที่มีการรับเงินหรือผลประโยชน์ตอบแทนคืนในระหว่างอายุกรมธรรม์:}
    \begin{itemize}
        \item กรณีได้รับเงินหรือผลประโยชน์ตอบแทนคืนทุกปี จะต้องไม่เกินร้อยละ 20 ของเบี้ยประกันชีวิตรายปี
        \item กรณีได้รับเงินหรือผลประโยชน์ตอบแทนคืนตามช่วงระยะเวลาที่ผู้รับประกันภัยกำหนดนอกจากข้อแรก เช่น ราย 2 ปี 3 ปี หรือ 5 ปี จะต้องไม่เกินร้อยละ 20 ของเบี้ยประกันชีวิตสะสมของแต่ละช่วงระยะเวลาที่ผู้รับประกันภัยกำหนดให้มีการจ่ายเงินหรือผลประโยชน์ตอบแทนคืน (ทั้งนี้ เงินหรือผลประโยชน์ตอบแทนคืนไม่รวมเงินปันผลตามกรมธรรม์ประกันชีวิตที่ได้รับตามเงื่อนไขในกรมธรรม์ประกันชีวิต หรือเงินหรือผลประโยชน์ตอบแทนที่จ่ายเมื่อสิ้นสุดการจ่ายเบี้ยประกันชีวิตแล้ว)
    \end{itemize}
    \item \textbf{กรณีผู้มีเงินได้มีคู่สมรสที่ไม่มีเงินได้:} นอกจากผู้มีเงินได้จะหักลดหย่อนเบี้ยประกันชีวิตในส่วนของตนตามที่จ่ายจริงได้ไม่เกิน 100,000 บาท (ตามมาตรา 47(1)(ง) ประกอบข้อ 2(61) แห่งกฎกระทรวงฉบับที่ 126) แล้ว ในกรณีที่คู่สมรสของผู้มีเงินได้มีการประกันชีวิตและความเป็นสามีภริยาได้มีอยู่ตลอดปีภาษี ผู้มีเงินได้มีสิทธิหักลดหย่อนสำหรับเบี้ยประกันชีวิตของคู่สมรสฝ่ายที่ไม่มีเงินได้ตามจำนวนที่จ่ายจริงแต่ไม่เกิน \hilight{10,000 บาท} ตามมาตรา 47(1)(ง) แห่งประมวลรัษฎากรเพิ่มเติมอีกด้วย (ได้เฉพาะส่วนตามมาตรา 47(1)(ง) ไม่ได้ส่วนเพิ่ม 90,000 บาทตามข้อ 2(61))
    \item \textbf{กรณีผู้มีเงินได้มีคู่สมรสที่มีเงินได้:}
    \begin{itemize}
        \item กรณีแยกยื่น --- ต่างฝ่ายต่างหักค่าลดหย่อนเบี้ยประกันชีวิตในส่วนของตนตามที่จ่ายจริงได้ไม่เกิน 100,000 บาท (ตามมาตรา 47(1)(ง) ประกอบข้อ 2(61) แห่งกฎกระทรวงฉบับที่ 126)
        \item กรณีรวมยื่น --- ผู้มีเงินได้หักลดหย่อนเบี้ยประกันชีวิตในส่วนของตนตามที่จ่ายจริงได้ไม่เกิน 100,000 บาท และสามารถหักค่าลดหย่อนเบี้ยประกันชีวิตในส่วนของคู่สมรสตามที่จ่ายจริงไม่เกิน 100,000 บาท (ตามมาตรา 47(1)(ง) ประกอบข้อ 2(61) แห่งกฎกระทรวงฉบับที่ 126) ได้ \hilight{(รวมหักได้สูงสุด 200,000 บาท)}
    \end{itemize}
\end{itemize}

\subsection*{ตัวอย่างการคำนวณ (สำหรับฝึกทำข้อสอบ)}

\textbf{ตัวอย่างที่ 1:} นาย ก. (ผู้มีเงินได้) มีเงินเดือน 2,000,000 บาท จ่ายเงินเข้ากองทุนสำรองเลี้ยงชีพ 200,000 บาท จ่ายเงินซื้อกองทุนรวมเพื่อการเลี้ยงชีพ (RMF) 200,000 บาท จ่ายเงินซื้อประกันชีวิตแบบบำนาญ 250,000 บาท คำถาม: นาย ก. จะหักลดหย่อน/ยกเว้นได้เป็นจำนวนเท่าใด

แนวคิดการวิเคราะห์: ต้องนำเพดานรวมตามข้อ 2(61) วรรคสองมาพิจารณาประกอบ คือเมื่อรวมเงินสะสมกองทุนสำรองเลี้ยงชีพ เงินค่าซื้อ RMF และเบี้ยประกันชีวิตแบบบำนาญแล้วต้องไม่เกิน 500,000 บาท ในกรณีนี้แม้นาย ก. จ่ายรวมกันทั้งสิ้น 650,000 บาท (200,000 + 200,000 + 250,000) ก็ตาม กฎหมายจำกัดเพดานรวมไว้เพียง 500,000 บาทเท่านั้น

\textbf{ตัวอย่างที่ 2:} นาง ข. (ผู้มีเงินได้) มีเงินเดือน 1,000,000 บาท จ่ายเงินเข้ากองทุนสำรองเลี้ยงชีพ 100,000 บาท จ่ายเงินซื้อกองทุนรวมเพื่อการเลี้ยงชีพ (RMF) 150,000 บาท จ่ายเงินซื้อประกันชีวิตแบบบำนาญ 280,000 บาท คำถาม: นาง ข. จะหักลดหย่อน/ยกเว้นได้เป็นจำนวนเท่าใด

แนวคิดการวิเคราะห์: ต้องพิจารณาเพดานของแต่ละประเภทก่อนแล้วจึงนำผลรวมมาเทียบกับเพดานรวม 500,000 บาท
\begin{itemize}
    \item RMF ไม่เกินร้อยละ 30 ของเงินได้พึงประเมิน = 300,000 บาท (เพดานตามสัดส่วน) แต่จ่ายจริงเพียง 150,000 บาท ดังนั้นหักได้ 150,000 บาทตามที่จ่ายจริง
    \item ประกันชีวิตแบบบำนาญ ไม่เกินร้อยละ 15 ของเงินได้พึงประเมิน = 150,000 บาท แต่จ่ายจริง 280,000 บาท ดังนั้นหักได้เพียง 150,000 บาทตามเพดานสัดส่วน
    \item กองทุนสำรองเลี้ยงชีพ จ่ายจริง 100,000 บาท หักได้ 100,000 บาท
    \item รวมทั้งสามรายการ = 100,000 + 150,000 + 150,000 = 400,000 บาท ซึ่งไม่เกินเพดานรวม 500,000 บาท จึงหักได้ทั้งหมด 400,000 บาท
\end{itemize}

\subsection*{ค่าลดหย่อนที่เกี่ยวเนื่อง}

\textbf{ค่าลดหย่อนเงินฝากที่มีประกันชีวิต (ข้อ 2(94))} เงินได้เท่าที่ผู้มีเงินได้จ่ายเป็นเงินฝากไว้กับธนาคารที่มีกฎหมายจัดตั้งขึ้นโดยเฉพาะในปีภาษีตามจำนวนที่จ่ายจริงแต่ไม่เกินหนึ่งแสนบาท ซึ่งการฝากเงินนั้นมีข้อตกลงว่าธนาคารผู้รับฝากเงินจะจ่ายเงินและผลประโยชน์ตามข้อตกลงโดยอาศัยความทรงชีพหรือมรณะของผู้ฝากเงิน และมีกำหนดระยะเวลาการฝากเงินตั้งแต่สิบปีขึ้นไป ทั้งนี้เมื่อรวมกับค่าลดหย่อนตามมาตรา 47(1)(ง) แห่งประมวลรัษฎากร หรือเงินได้ตาม (61) วรรคหนึ่งแล้ว ต้องไม่เกิน \hilight{100,000 บาท}

\textbf{ค่าลดหย่อนเบี้ยประกันสุขภาพ (ข้อ 2(97))} เงินได้เท่าที่ผู้มีเงินได้จ่ายเป็นเบี้ยประกันภัยในปีภาษีให้แก่บริษัทประกันชีวิตหรือบริษัทประกันวินาศภัยที่ประกอบกิจการในราชอาณาจักร สำหรับการประกันสุขภาพของผู้มีเงินได้ ตามจำนวนที่จ่ายจริงแต่ไม่เกินสองหมื่นห้าพันบาท ซึ่งเมื่อรวมกับค่าลดหย่อนตามมาตรา 47(1)(ง) แห่งประมวลรัษฎากร หรือเงินได้ที่ได้รับยกเว้นตาม (61) วรรคหนึ่ง หรือ (94) แล้วแต่กรณี ต้องไม่เกินหนึ่งแสนบาท ทั้งนี้สำหรับเบี้ยประกันภัยที่ได้จ่ายตั้งแต่วันที่ 1 มกราคม พ.ศ. 2563 เป็นต้นไป (แก้ไขเพิ่มเติมโดยกฎกระทรวง ฉบับที่ 365 พ.ศ. 2563)


\newpage
\section{ค่าลดหย่อนเงินสะสมเข้ากองทุนสำรองเลี้ยงชีพ}

ฐานทางกฎหมายของค่าลดหย่อนนี้มีสองชั้นเช่นเดียวกับเบี้ยประกันชีวิต คือมาตรา 47(1)(ช) แห่งประมวลรัษฎากร และข้อ 2(35) แห่งกฎกระทรวงฉบับที่ 126

\subsection*{ฐานทางกฎหมายที่ 1: มาตรา 47(1)(ช)}

\begin{quote}
{\itshape "(ช) เงินสะสมที่จ่ายเข้ากองทุนสำรองเลี้ยงชีพ ซึ่งเป็นไปตามหลักเกณฑ์ วิธีการ และเงื่อนไขที่กำหนดโดยกฎกระทรวงตามมาตรา 65 ตรี (2) ตามจำนวนที่จ่ายจริงแต่ไม่เกิน 10,000 บาท (แก้ไขโดยประกาศคณะรักษาความสงบเรียบร้อยแห่งชาติ ฉบับที่ 37 พ.ศ. 2534 ใช้บังคับ 1 ม.ค. 2534 เป็นต้นไป)

ในกรณีสามีหรือภริยาของผู้มีเงินได้จ่ายเงินสะสมเข้ากองทุนสำรองเลี้ยงชีพตามวรรคหนึ่ง และความเป็นสามีภริยาได้มีอยู่ตลอดปีภาษี ให้หักลดหย่อนได้ด้วยสำหรับเงินสะสมของสามีหรือภริยาที่จ่ายเข้ากองทุนสำรองเลี้ยงชีพนั้นตามเกณฑ์ในวรรคหนึ่ง"} (พระราชกำหนดแก้ไขเพิ่มเติม ฉบับที่ 12 พ.ศ. 2526 ใช้บังคับปีภาษี 2527 เป็นต้นไป)
\end{quote}

\subsection*{ฐานทางกฎหมายที่ 2: ข้อ 2(35) แห่งกฎกระทรวงฉบับที่ 126}

\begin{quote}
\textit{"(35) เงินได้เท่าที่ลูกจ้างจ่ายเป็นเงินสะสมเข้ากองทุนสำรองเลี้ยงชีพ ตามกฎหมายว่าด้วยกองทุนสำรองเลี้ยงชีพในอัตราไม่เกินร้อยละสิบห้าของค่าจ้าง เฉพาะส่วนที่เกินหนึ่งหมื่นบาท แต่ไม่เกินสี่แสนเก้าหมื่นบาท สำหรับปีภาษีนั้น ทั้งนี้สำหรับเงินได้พึงประเมินที่ได้รับตั้งแต่วันที่ 1 มกราคม พ.ศ. 2551 เป็นต้นไป"}
\end{quote}

กล่าวคือ เมื่อรวมฐานทั้งสองแล้ว ผู้มีเงินได้สามารถหักค่าเงินสะสมเข้ากองทุนสำรองเลี้ยงชีพได้สูงสุด \hilight{ไม่เกิน 500,000 บาท} (10,000 บาทตามมาตรา 47(1)(ช) บวกส่วนเพิ่ม 490,000 บาทตามข้อ 2(35))

\subsection*{หมายเหตุประกอบ: กองทุนประเภทอื่นที่มีเพดานเกี่ยวข้อง}

\textbf{กองทุนบำเหน็จบำนาญข้าราชการ (ข้อ 2(43))} เงินได้เท่าที่สมาชิกกองทุนบำเหน็จบำนาญข้าราชการจ่ายเป็นเงินสะสมเข้ากองทุนบำเหน็จบำนาญข้าราชการตามกฎหมายว่าด้วยกองทุนบำเหน็จบำนาญข้าราชการ เฉพาะส่วนที่ไม่เกินห้าแสนบาท สำหรับปีภาษีนั้น ทั้งนี้สำหรับเงินได้พึงประเมินที่ได้รับตั้งแต่วันที่ 1 มกราคม พ.ศ. 2551 เป็นต้นไป (แก้ไขเพิ่มเติมโดยกฎกระทรวงฉบับที่ 266 พ.ศ. 2551)

\textbf{กองทุนสงเคราะห์โรงเรียนเอกชน (ข้อ 2(54))} เงินได้เท่าที่ผู้อำนวยการ ผู้บริหาร ครู หรือบุคลากรทางการศึกษาในโรงเรียนเอกชน จ่ายเป็นเงินสะสมเข้ากองทุนสงเคราะห์ตามกฎหมายว่าด้วยโรงเรียนเอกชน เฉพาะส่วนที่ไม่เกินห้าแสนบาท สำหรับปีภาษีนั้น (แก้ไขเพิ่มเติมโดยกฎกระทรวงฉบับที่ 266 พ.ศ. 2551)

\textbf{กองทุนการออมแห่งชาติ (ข้อ 2(90))} เงินได้เท่าที่สมาชิกกองทุนการออมแห่งชาติจ่ายเป็นเงินสะสมเข้ากองทุนการออมแห่งชาติ ตามกฎหมายว่าด้วยกองทุนการออมแห่งชาติ ตามจำนวนที่จ่ายจริงแต่ไม่เกินห้าแสนบาท สำหรับปีภาษีนั้น เงินได้ที่ได้รับยกเว้นตามวรรคหนึ่ง เมื่อรวมกับเงินได้ที่ได้รับยกเว้นไม่ต้องรวมคำนวณเพื่อเสียภาษีเงินได้สำหรับกรณีที่ผู้มีเงินได้จ่ายเป็นเงินสะสมเข้ากองทุนสำรองเลี้ยงชีพ เงินสะสมเข้ากองทุนบำเหน็จบำนาญข้าราชการ เงินสะสมเข้ากองทุนสงเคราะห์โรงเรียนเอกชน หรือเงินค่าซื้อหน่วยลงทุนในกองทุนรวมเพื่อการเลี้ยงชีพ หรือเบี้ยประกันภัยสำหรับการประกันชีวิตแบบบำนาญแล้ว ต้องไม่เกินห้าแสนบาทในปีภาษีเดียวกัน (ดูประกอบประกาศอธิบดีกรมสรรพากรเกี่ยวกับภาษีเงินได้ ฉบับที่ 274)

\subsection*{ค่าลดหย่อนการซื้อหน่วยลงทุนในกองทุนต่าง ๆ}

\textbf{RMF --- กองทุนรวมเพื่อการเลี้ยงชีพ (ข้อ 2(55))} เงินได้เท่าที่จ่ายเป็นค่าซื้อหน่วยลงทุนในกองทุนรวมเพื่อการเลี้ยงชีพตามกฎหมายว่าด้วยหลักทรัพย์และตลาดหลักทรัพย์ ในอัตราไม่เกินร้อยละสามสิบของเงินได้พึงประเมิน เฉพาะส่วนที่ไม่เกินห้าแสนบาท สำหรับปีภาษีนั้น โดยผู้มีเงินได้ต้องถือหน่วยลงทุนดังกล่าวมาแล้วไม่น้อยกว่าห้าปีนับตั้งแต่วันซื้อหน่วยลงทุนครั้งแรกและไถ่ถอนหน่วยลงทุนนั้น เมื่อผู้มีเงินได้มีอายุไม่ต่ำกว่าห้าสิบห้าปีบริบูรณ์ ทั้งนี้สำหรับเงินได้พึงประเมินที่ได้รับตั้งแต่วันที่ 1 มกราคม พ.ศ. 2563 เป็นต้นไป (แก้ไขเพิ่มเติมโดยกฎกระทรวงฉบับที่ 357 พ.ศ. 2563) (ดูประกอบประกาศอธิบดีกรมสรรพากรเกี่ยวกับภาษีเงินได้ ฉบับที่ 401) ในกรณีที่ผู้มีเงินได้ถือหน่วยลงทุนไม่ครบห้าปีนับตั้งแต่วันที่ซื้อหน่วยลงทุนครั้งแรก หรือไถ่ถอนหน่วยลงทุนก่อนที่ผู้มีเงินได้มีอายุครบห้าสิบห้าปีบริบูรณ์ ให้ผู้มีเงินได้นั้นหมดสิทธิได้รับยกเว้นและต้องเสียภาษีเงินได้สำหรับเงินได้ที่ได้รับยกเว้นภาษีไปแล้วด้วย

\textbf{LTF --- กองทุนรวมหุ้นระยะยาว (ข้อ 2(66))} เงินได้เท่าที่จ่ายเป็นค่าซื้อหน่วยลงทุนในกองทุนรวมหุ้นระยะยาวตามกฎหมายว่าด้วยหลักทรัพย์และตลาดหลักทรัพย์ ในอัตราไม่เกินร้อยละสิบห้าของเงินได้พึงประเมิน เฉพาะส่วนที่ไม่เกินห้าแสนบาท สำหรับปีภาษีนั้น โดยผู้มีเงินได้ต้องถือหน่วยลงทุนในกองทุนรวมหุ้นระยะยาวไม่น้อยกว่าเจ็ดปีปฏิทิน ทั้งนี้สำหรับเงินได้พึงประเมินที่ได้รับตั้งแต่วันที่ 1 มกราคม พ.ศ. 2559 ถึงวันที่ 31 ธันวาคม พ.ศ. 2562 \hilight{ปัจจุบัน LTF ไม่มีสิทธิประโยชน์ใหม่อีกแล้ว} สิทธิเดิมตามข้อ 2(66) ยังคงใช้บังคับต่อไปเฉพาะเงินได้เท่าที่จ่ายเป็นค่าซื้อหน่วยลงทุนในกองทุนรวมหุ้นระยะยาวที่ซื้อก่อนวันที่ 1 มกราคม พ.ศ. 2563 (ดูประกอบประกาศอธิบดีกรมสรรพากรเกี่ยวกับภาษีเงินได้ ฉบับที่ 276)

\textbf{SSF --- กองทุนรวมเพื่อการออม (ข้อ 2(102))} เงินได้เท่าที่จ่ายเป็นค่าซื้อหน่วยลงทุนในกองทุนรวมเพื่อการออมตามกฎหมายว่าด้วยหลักทรัพย์และตลาดหลักทรัพย์ ในอัตราไม่เกินร้อยละสามสิบของเงินได้พึงประเมิน เฉพาะส่วนที่ไม่เกินสองแสนบาท สำหรับปีภาษีนั้น โดยผู้มีเงินได้ต้องถือหน่วยลงทุนในกองทุนรวมเพื่อการออมไม่น้อยกว่าสิบปีนับตั้งแต่วันที่ซื้อหน่วยลงทุน ทั้งนี้สำหรับเงินได้พึงประเมินที่ได้รับตั้งแต่วันที่ 1 มกราคม พ.ศ. 2563 ถึงวันที่ 31 ธันวาคม พ.ศ. 2567 (ดูประกอบประกาศอธิบดีกรมสรรพากรเกี่ยวกับภาษีเงินได้ ฉบับที่ 369)

\textbf{Thai ESG / TESG --- กองทุนรวมไทยเพื่อความยั่งยืน (ข้อ 2(105))} เงินได้เท่าที่จ่ายเป็นค่าซื้อหน่วยลงทุนในกองทุนรวมไทยเพื่อความยั่งยืนตามกฎหมายว่าด้วยหลักทรัพย์และตลาดหลักทรัพย์ ในอัตราไม่เกินร้อยละสามสิบของเงินได้พึงประเมิน เฉพาะส่วนที่ไม่เกินหนึ่งแสนบาท สำหรับปีภาษีนั้น โดยผู้มีเงินได้ต้องถือหน่วยลงทุนไม่น้อยกว่าแปดปีนับตั้งแต่วันที่ซื้อหน่วยลงทุน ทั้งนี้สำหรับเงินได้พึงประเมินที่ได้รับตั้งแต่วันที่ 21 พฤศจิกายน พ.ศ. 2566 ถึงวันที่ 31 ธันวาคม พ.ศ. 2575 และในปีภาษี 2567 ถึงปีภาษี 2569 หากผู้มีเงินได้มีการซื้อหน่วยลงทุนระหว่างวันที่ 1 มกราคม พ.ศ. 2567 ถึงวันที่ 31 ธันวาคม พ.ศ. 2569 ให้ได้รับยกเว้นในอัตราไม่เกินร้อยละสามสิบของเงินได้พึงประเมิน เฉพาะส่วนที่ไม่เกินสามแสนบาท โดยต้องถือหน่วยลงทุนไม่น้อยกว่าห้าปี (แก้ไขเพิ่มเติมโดยกฎกระทรวงฉบับที่ 395 พ.ศ. 2567) (ดูประกอบประกาศอธิบดีกรมสรรพากรเกี่ยวกับภาษีเงินได้ ฉบับที่ 442)

\textbf{Thai ESGX / TESGX (ข้อ 2(107))} กรณีซื้อหน่วยลงทุนใหม่ ผู้มีเงินได้ที่ซื้อหน่วยลงทุนในกองทุนรวมไทยเพื่อความยั่งยืนแบบพิเศษ ได้รับยกเว้นภาษีในอัตราไม่เกินร้อยละสามสิบของเงินได้พึงประเมิน เฉพาะส่วนที่ไม่เกินสามแสนบาท สำหรับการซื้อภายในสองเดือนนับแต่วันเปิดขายหน่วยลงทุนครั้งแรก แต่ไม่เกินวันที่ 30 มิถุนายน พ.ศ. 2568 และกรณีสับเปลี่ยนหน่วยลงทุน LTF เดิมเป็น Thai ESGX ผู้มีเงินได้จะได้รับยกเว้นภาษีเงินได้เท่ากับมูลค่าหน่วยลงทุนที่สับเปลี่ยนแต่ไม่เกินห้าแสนบาท โดยแบ่งสิทธิยกเว้นเป็นช่วงเวลา คือ ปีภาษี 2568 ยกเว้นเฉพาะส่วนที่ไม่เกินสามแสนบาท และปีภาษี 2569 ถึงปีภาษี 2572 ยกเว้นเฉพาะส่วนที่เกินสามแสนบาทแต่ไม่เกินห้าแสนบาท โดยให้ได้รับยกเว้นเป็นจำนวนเท่า ๆ กันในแต่ละปีภาษี (ดูประกอบประกาศอธิบดีกรมสรรพากรเกี่ยวกับภาษีเงินได้ ฉบับที่ 442)

\subsection*{ตารางเปรียบเทียบกองทุนเพื่อการลงทุนระยะยาวที่ได้รับสิทธิประโยชน์ทางภาษี}

\begin{longtable}{|p{2.3cm}|p{2.7cm}|p{2.7cm}|p{2.7cm}|p{3.3cm}|}
\hline
\textbf{รายการ} & \textbf{RMF} & \textbf{SSF} & \textbf{Thai ESG} & \textbf{Thai ESGX} \\
\hline
\endhead
สัดส่วนสูงสุดของการลดหย่อน & ไม่เกิน 30\% ของเงินได้ฯ & ไม่เกิน 30\% ของเงินได้ฯ & ไม่เกิน 30\% ของเงินได้ฯ & เงินใหม่ไม่เกิน 30\% ของเงินได้ฯ \\
\hline
จำนวนสูงสุดที่ลดหย่อนได้ (บาท) & 500,000 (รวมกับกองทุนเพื่อการเกษียณอื่น ๆ ต้องไม่เกิน 500,000) & 200,000 (รวมกับกองทุนเพื่อการเกษียณอื่น ๆ ต้องไม่เกิน 500,000) & 100,000 (300,000 ช่วงปี 2567--2569) & เงินใหม่ 300,000; LTF เดิมที่สับเปลี่ยน 500,000 (ปี 2568 = 300,000 ส่วนที่เหลือทยอยหักเท่า ๆ กันปี 2569--2572) \\
\hline
ระยะเวลาการลงทุน & ไม่น้อยกว่า 5 ปี และขายได้เมื่ออายุครบ 55 ปีบริบูรณ์ & ไม่น้อยกว่า 10 ปี & ไม่น้อยกว่า 8 ปี (ไม่น้อยกว่า 5 ปี ในกรณีพิเศษ) & ไม่น้อยกว่า 5 ปี \\
\hline
เงื่อนไขการซื้อ & ไม่กำหนดเงินลงทุนขั้นต่ำ แต่ต้องซื้อต่อเนื่องทุกปีหรือปีเว้นปี & ไม่กำหนดเงินลงทุนขั้นต่ำ & ไม่กำหนดเงินลงทุนขั้นต่ำ & ไม่กำหนดเงินลงทุนขั้นต่ำ แต่ต้องซื้อหรือสับเปลี่ยนภายในวันที่ 30 มิ.ย. 2568 \\
\hline
สินทรัพย์ที่ลงทุน & สินทรัพย์ทุกประเภท & สินทรัพย์ทุกประเภท & หุ้นและตราสารหนี้ที่ได้ SET ESG Rating มากกว่า 80\% & หุ้นและตราสารหนี้ที่ได้ SET ESG Rating มากกว่า 80\% โดยเป็นหุ้นมากกว่า 65\% \\
\hline
ปีที่สามารถใช้สิทธิ & 2563 เป็นต้นไป & 2563--2567 & 2566, 2570--2575 (หรือ 2567--2569) & 2568--2572 \\
\hline
\end{longtable}

\newpage
\section{ค่าลดหย่อนดอกเบี้ยเงินกู้ยืมเพื่อที่อยู่อาศัย}

ฐานทางกฎหมาย: มาตรา 47(1)(ซ) แห่งประมวลรัษฎากร และข้อ 2(53) แห่งกฎกระทรวงฉบับที่ 126

\subsection*{มาตรา 47(1)(ซ)}

\begin{quote}
\textit{"(ซ) ดอกเบี้ยเงินกู้ยืมที่ผู้มีเงินได้จ่ายให้แก่ธนาคารหรือสถาบันการเงินอื่น บริษัทประกันชีวิต สหกรณ์หรือนายจ้าง สำหรับการกู้ยืมเงินเพื่อซื้อ เช่าซื้อ หรือสร้างอาคารอยู่อาศัย โดยจำนองอาคารที่ซื้อหรือสร้างเป็นประกันการกู้ยืมนั้น ตามจำนวนที่จ่ายจริงแต่ไม่เกิน 10,000 บาท ทั้งนี้ตามหลักเกณฑ์และวิธีการที่อธิบดีกำหนดโดยอนุมัติรัฐมนตรีและประกาศในราชกิจจานุเบกษา อาคารดังกล่าวให้หมายความรวมถึงอาคารพร้อมที่ดินด้วย"} (แก้ไขโดยประกาศคณะรักษาความสงบเรียบร้อยแห่งชาติ ฉบับที่ 37 พ.ศ. 2534 ใช้บังคับ 1 ม.ค. 2534 เป็นต้นไป)
\end{quote}

\subsection*{ข้อ 2(53) แห่งกฎกระทรวงฉบับที่ 126}

\begin{quote}
{\itshape "(53) เงินได้เท่าที่ได้จ่ายเป็นดอกเบี้ยเงินกู้ยืมให้แก่ธนาคารหรือสถาบันการเงินอื่น บริษัทประกันชีวิต สหกรณ์ หรือนายจ้าง สำหรับการกู้ยืมเงินเพื่อซื้อ เช่าซื้อ หรือสร้างอาคารที่อยู่อาศัย โดยจำนองอาคารที่ซื้อหรือสร้างเป็นประกันการกู้ยืมนั้น ตามจำนวนที่จ่ายจริงในส่วนที่เกินหนึ่งหมื่นบาทแต่ไม่เกินเก้าหมื่นบาท และเฉพาะดอกเบี้ยเงินกู้ยืมที่ได้จ่ายตั้งแต่วันที่ 1 มกราคม พ.ศ. 2550 เป็นต้นไป ทั้งนี้ตามหลักเกณฑ์และวิธีการที่อธิบดีกำหนด

กรณีที่ผู้มีเงินได้หักลดหย่อนตามมาตรา 47(1)(ซ) แห่งประมวลรัษฎากร หรือได้รับยกเว้นไม่ต้องนำเงินได้ตาม (52) หรือ (59) รวมคำนวณเพื่อเสียภาษีเงินได้ เงินได้ที่ได้รับยกเว้นตามวรรคหนึ่งเมื่อรวมกับค่าลดหย่อนตามมาตรา 47(1)(ซ) แห่งประมวลรัษฎากร หรือเงินได้ตาม (52) หรือ (59) แล้วแต่กรณี ต้องไม่เกินหนึ่งแสนบาท

อาคารตามวรรคหนึ่งให้หมายความรวมถึงอาคารพร้อมที่ดินด้วย"} (แก้ไขเพิ่มเติมโดยกฎกระทรวงฉบับที่ 264 พ.ศ. 2550 ตั้งแต่ 1 มกราคม 2550 เป็นต้นไป)
\end{quote}

กล่าวคือ เมื่อรวมฐานทั้งสองแล้ว ผู้มีเงินได้หักดอกเบี้ยเงินกู้ยืมเพื่อที่อยู่อาศัยได้สูงสุด \hilight{ไม่เกิน 100,000 บาท} (10,000 บาทตามมาตรา 47(1)(ซ) บวกส่วนเพิ่ม 90,000 บาทตามข้อ 2(53)) (ดูประกาศกรมสรรพากร เรื่อง หลักฐานการหักลดหย่อนสำหรับดอกเบี้ยเงินกู้ยืม)

\subsection*{หมายเหตุประกอบ}

\textbf{ดอกเบี้ยเงินกู้ยืมกองทุนรวมอสังหาริมทรัพย์เพื่อที่อยู่อาศัย (ข้อ 2(52))} เงินได้เท่าที่ได้จ่ายเป็นดอกเบี้ยเงินกู้ยืม สำหรับการกู้ยืมเงินเพื่อซื้อ เช่าซื้อ หรือสร้างอาคารที่อยู่อาศัย โดยจำนองอาคารที่ซื้อหรือสร้างเป็นประกันการกู้ยืมนั้น ตามจำนวนที่จ่ายจริงแต่ไม่เกินหนึ่งแสนบาท ตามหลักเกณฑ์และวิธีการที่อธิบดีกำหนด ทั้งนี้ เฉพาะดอกเบี้ยเงินกู้ยืมที่ได้จ่ายตั้งแต่วันที่ 1 มกราคม พ.ศ. 2550 เป็นต้นไป ซึ่งจ่ายให้แก่ (ก) กองทุนรวมอสังหาริมทรัพย์เพื่อแก้ไขปัญหาในระบบสถาบันการเงินที่จัดตั้งขึ้นตามกฎหมายว่าด้วยหลักทรัพย์และตลาดหลักทรัพย์ (ข) กองทุนรวมเพื่อแก้ไขปัญหาในระบบสถาบันการเงินที่จัดตั้งขึ้นตามกฎหมายว่าด้วยหลักทรัพย์และตลาดหลักทรัพย์ หรือ (ค) นิติบุคคลเฉพาะกิจซึ่งจัดตั้งขึ้นเพื่อดำเนินการแปลงสินทรัพย์เป็นหลักทรัพย์ตามกฎหมายว่าด้วยนิติบุคคลเฉพาะกิจเพื่อการแปลงสินทรัพย์เป็นหลักทรัพย์ เฉพาะกรณีที่นิติบุคคลเฉพาะกิจดังกล่าวเข้ารับช่วงสิทธิเป็นเจ้าหนี้เงินกู้แทนกองทุนรวมตาม (ก) หรือ (ข) ธนาคารหรือสถาบันการเงินอื่น บริษัทประกันชีวิต สหกรณ์ หรือนายจ้าง เมื่อรวมกับค่าลดหย่อนตามมาตรา 47(1)(ซ) หรือเงินได้ตาม (53) หรือ (59) แล้วแต่กรณี ต้องไม่เกินหนึ่งแสนบาท (แก้ไขเพิ่มเติมโดยกฎกระทรวงฉบับที่ 264 พ.ศ. 2550)

\textbf{ดอกเบี้ยเงินกู้ยืมกองทุนบำเหน็จบำนาญข้าราชการเพื่อที่อยู่อาศัย (ข้อ 2(59))} เงินได้เท่าที่ได้จ่ายเป็นดอกเบี้ยเงินกู้ยืมให้แก่กองทุนบำเหน็จบำนาญข้าราชการตามกฎหมายว่าด้วยกองทุนบำเหน็จบำนาญข้าราชการ สำหรับการกู้ยืมเงินเพื่อซื้อ เช่าซื้อ หรือสร้างอาคารที่อยู่อาศัย ตามจำนวนที่จ่ายจริงแต่ไม่เกินหนึ่งแสนบาท และเฉพาะดอกเบี้ยเงินกู้ยืมที่ได้จ่ายตั้งแต่วันที่ 1 มกราคม พ.ศ. 2550 เป็นต้นไป ทั้งนี้ตามหลักเกณฑ์และวิธีการที่อธิบดีกำหนด เมื่อรวมกับค่าลดหย่อนตามมาตรา 47(1)(ซ) หรือเงินได้ตาม (52) หรือ (53) แล้วแต่กรณี ต้องไม่เกินหนึ่งแสนบาท (แก้ไขเพิ่มเติมโดยกฎกระทรวงฉบับที่ 264 พ.ศ. 2550)

\hilight{ข้อสำคัญสำหรับข้อสอบ:} ดอกเบี้ยเงินกู้ยืมเพื่อที่อยู่อาศัยทั้งสามฐาน (มาตรา 47(1)(ซ)/ข้อ 2(53), ข้อ 2(52), และข้อ 2(59)) มีเพดานรวมกันไม่เกิน 100,000 บาทเสมอ แม้ผู้มีเงินได้จะมีการกู้ยืมจากหลายแหล่งพร้อมกัน

\newpage
\section{ค่าลดหย่อนเงินสมทบที่จ่ายเข้ากองทุนประกันสังคม}

ฐานทางกฎหมาย: มาตรา 47(1)(ฌ) แห่งประมวลรัษฎากร

หักลดหย่อนให้สำหรับผู้มีเงินได้ตามจำนวนที่จ่ายจริง (เงินสมทบของลูกจ้างผู้ประกันตน) โดยไม่มีเพดานกำหนดไว้ตรง ๆ แต่ในทางปฏิบัติจำนวนเงินสมทบถูกจำกัดด้วยอัตราและฐานเงินเดือนตามพระราชบัญญัติประกันสังคม

มาตรา 47(1)(ฌ) บัญญัติว่า

\begin{quote}
\textit{"(ฌ) เงินสมทบที่ผู้ประกันตนจ่ายเข้ากองทุนประกันสังคมตามกฎหมายว่าด้วยการประกันสังคมตามจำนวนที่จ่ายจริง"} (พระราชบัญญัติแก้ไขเพิ่มเติม ฉบับที่ 27 พ.ศ. 2533 ใช้บังคับ 3 ก.ย. 2533 เป็นต้นไป)

\textit{"ในกรณีสามีหรือภริยาของผู้มีเงินได้ ซึ่งเป็นผู้ประกันตนจ่ายเงินสมทบเข้ากองทุนประกันสังคมตามวรรคหนึ่ง และความเป็นสามีภริยาได้มีอยู่ตลอดปีภาษี ให้หักลดหย่อนได้ด้วยสำหรับเงินสมทบของสามีหรือภริยาที่จ่ายเข้ากองทุนประกันสังคมดังกล่าว ตามเกณฑ์ในวรรคหนึ่ง"} (พระราชกำหนดแก้ไขเพิ่มเติม ฉบับที่ 16 พ.ศ. 2534 ใช้บังคับปีภาษี 2535 เป็นต้นไป)
\end{quote}

ดูพระราชบัญญัติประกันสังคม พ.ศ. 2533 และกฎกระทรวงที่เกี่ยวข้อง โดยอัตราเงินสมทบปกติคือร้อยละ 5 ของฐานเงินเดือน 15,000 บาท (คิดเป็น 750 บาทต่อเดือน) ทั้งนี้ในบางช่วงมีการปรับลดอัตราเงินสมทบเป็นการชั่วคราวเนื่องจากสถานการณ์พิเศษ เช่น ร้อยละ 1 (มี.ค.--พ.ค. 2563) ร้อยละ 2 (ก.ย.--พ.ย. 2563) ร้อยละ 3 (ม.ค. 2564) ร้อยละ 0.5 (ก.พ.--มี.ค. 2564) และร้อยละ 2.5 (มิ.ย.--ส.ค. 2564) (ดูประกาศอธิบดีกรมสรรพากรเกี่ยวกับภาษีเงินได้ ฉบับที่ 382)

\newpage
\section{ค่าลดหย่อนการอุปการะเลี้ยงดูบิดามารดา}

ฐานทางกฎหมาย: มาตรา 47(1)(ญ) แห่งประมวลรัษฎากร

หักลดหย่อนให้สำหรับบิดามารดาของผู้มีเงินได้ และบิดามารดาของคู่สมรส คนละ \hilight{30,000 บาท} (สูงสุด 4 คน เท่ากับ 120,000 บาท)

มาตรา 47(1)(ญ) บัญญัติว่า

\begin{quote}
\textit{"(ญ) ค่าอุปการะเลี้ยงดูบิดามารดาของผู้มีเงินได้ รวมทั้งบิดามารดาของสามีหรือภริยาของผู้มีเงินได้คนละสามหมื่นบาท โดยบุคคลดังกล่าวต้องมีอายุหกสิบปีขึ้นไป มีรายได้ไม่เพียงพอแก่การยังชีพ และอยู่ในความอุปการะเลี้ยงดูของผู้มีเงินได้ ทั้งนี้ตามหลักเกณฑ์ วิธีการ และเงื่อนไขที่อธิบดีประกาศกำหนด ค่าอุปการะเลี้ยงดูตามวรรคหนึ่ง ให้หักลดหย่อนสำหรับเงินได้พึงประเมินประจำปี พ.ศ. 2547 ที่จะต้องยื่นรายการในปี พ.ศ. 2548 เป็นต้นไป"} (พระราชบัญญัติแก้ไขเพิ่มเติม ฉบับที่ 36 พ.ศ. 2548 ใช้บังคับ 13 มกราคม 2548 เป็นต้นไป) (ดูประกาศอธิบดีกรมสรรพากรเกี่ยวกับภาษีเงินได้ ฉบับที่ 136) (ดูประกาศกรมสรรพากร เรื่อง หลักฐานการหักลดหย่อนค่าอุปการะเลี้ยงดูบิดามารดาของผู้มีเงินได้รวมทั้งบิดามารดาของสามีหรือภริยาของผู้มีเงินได้)
\end{quote}

\subsection*{สรุปเงื่อนไขต่าง ๆ ที่สำคัญกรณีค่าลดหย่อนอุปการะเลี้ยงดูบิดามารดา}

\begin{enumerate}
    \item \textbf{ค่าลดหย่อน:} ให้หักลดหย่อน "บิดาของผู้มีเงินได้" "มารดาของผู้มีเงินได้" "บิดาของคู่สมรสของผู้มีเงินได้" และ "มารดาของคู่สมรสของผู้มีเงินได้" ได้คนละ 30,000 บาท
    \item \textbf{อายุ:} บิดามารดาฯ ต้องมีอายุไม่ต่ำกว่า 60 ปี
    \item \textbf{ฐานะ:} เงินได้พึงประเมินซึ่งไม่ได้รับการยกเว้นภาษีที่บิดามารดาฯ แต่ละคนได้รับในแต่ละปีภาษี ต้องไม่เกิน 30,000 บาท
    \item \textbf{บุตรชอบด้วยกฎหมาย:} ผู้มีเงินได้หรือสามีหรือภริยาของผู้มีเงินได้จะต้องเป็นบุตรชอบด้วยกฎหมายของบิดามารดาที่ผู้มีเงินได้ใช้สิทธิหักลดหย่อน (บุตรบุญธรรมไม่มีสิทธิหักค่าลดหย่อนนี้)
    \item \textbf{หลักฐานรับรองการอุปการะเลี้ยงดู (ล.ย.03):} กรณีผู้มีเงินได้หลายคนอุปการะเลี้ยงดูบิดามารดาคนเดียวกัน ให้ผู้มีเงินได้คนใดคนหนึ่งเพียงคนเดียวที่มีหลักฐานรับรองการอุปการะเลี้ยงดูจากบิดามารดาดังกล่าวเป็นผู้มีสิทธิหักลดหย่อนค่าอุปการะเลี้ยงดูบิดามารดาดังกล่าวนั้นในแต่ละปีภาษี
    \item \textbf{บิดามารดาของคู่สมรส:} สามารถนำมาลดหย่อนได้กรณีที่คู่สมรสไม่มีรายได้ หรือกรณีที่คู่สมรสมีรายได้แต่ไม่ได้ใช้สิทธิแยกยื่นรายการและเสียภาษีต่างหาก (กรณีรวมยื่น) ทั้งนี้หากคู่สมรสมีรายได้และใช้สิทธิแยกยื่นรายการและเสียภาษีต่างหากฯ ให้สามีและภริยาต่างฝ่ายต่างหักลดหย่อนค่าอุปการะเลี้ยงดูบิดามารดาของตนคนละ 30,000 บาท
    \item \textbf{กรณีผู้มีเงินได้มิได้เป็นผู้อยู่ในประเทศไทย:} ให้หักลดหย่อนค่าอุปการะเลี้ยงดูบิดามารดาได้เฉพาะบิดามารดาที่อยู่ในประเทศไทย
\end{enumerate}

\subsection*{ตัวอย่างประกอบความเข้าใจ}

\textbf{ตัวอย่างที่ 1:} สามีเป็นผู้มีเงินได้และไม่ใช่ผู้อยู่ในประเทศไทย ภริยาไม่มีเงินได้แต่เป็นผู้อยู่ในประเทศไทย บุตรไม่มีเงินได้และไม่ใช่ผู้อยู่ในประเทศไทย บิดามารดาของภริยาเป็นผู้อยู่ในประเทศไทย บิดามารดาของสามีไม่ใช่ผู้อยู่ในประเทศไทย

ผลลัพธ์: เข้ากรณีตามมาตรา 47(3) สามารถหักค่าลดหย่อนคู่สมรสและบุตรได้เฉพาะกรณีที่คู่สมรส/บุตรเป็นผู้อยู่ในประเทศไทยเท่านั้น ดังนั้นสามีหักได้ ส่วนตัว 60,000 บาท คู่สมรส 60,000 บาท (เพราะภริยาอยู่ในประเทศไทย) บิดามารดาคู่สมรส 30,000+30,000 บาท (เพราะอยู่ในประเทศไทย) แต่หักบุตรไม่ได้และหักบิดามารดาของตนเองไม่ได้ (เพราะทั้งคู่ไม่ได้อยู่ในประเทศไทย)

\textbf{ตัวอย่างที่ 2:} สามีเป็นผู้มีเงินได้และเป็นผู้อยู่ในประเทศไทย (สมมติฐานอื่นเหมือนตัวอย่างที่ 1)

ผลลัพธ์: ไม่เข้ากรณีตามมาตรา 47(3) จึงหักค่าลดหย่อนคู่สมรสและบุตรได้ทั้งกรณีที่คู่สมรสและบุตรเป็นผู้อยู่ในประเทศไทย และกรณีที่คู่สมรสและบุตรไม่ใช่ผู้อยู่ในประเทศไทย สามีหักได้ ส่วนตัว 60,000 บาท คู่สมรส 60,000 บาท บุตร 30,000 บาท บิดามารดาผู้มีเงินได้ 30,000+30,000 บาท และบิดามารดาคู่สมรส 30,000+30,000 บาท (หักได้ทุกรายการเต็มจำนวน)

\newpage
\section{ค่าลดหย่อนเบี้ยประกันสุขภาพบิดามารดา}

ฐานทางกฎหมาย: ข้อ 2(76) แห่งกฎกระทรวงฉบับที่ 126

\begin{quote}
\textit{"(76) เงินได้เท่าที่ผู้มีเงินได้จ่ายเป็นเบี้ยประกันภัยให้แก่บริษัทประกันชีวิตหรือบริษัทประกันวินาศภัยที่ประกอบกิจการในราชอาณาจักรตามจำนวนที่จ่ายจริงแต่ไม่เกินหนึ่งหมื่นห้าพันบาท สำหรับการประกันสุขภาพบิดามารดาของผู้มีเงินได้ รวมทั้งบิดามารดาของสามีหรือภริยาของผู้มีเงินได้ซึ่งมีรายได้ไม่เพียงพอแก่การยังชีพ ทั้งนี้ต้องเป็นเบี้ยประกันภัยที่ได้จ่ายในปี พ.ศ. 2549 เป็นต้นไป และให้เป็นไปตามหลักเกณฑ์และวิธีการที่อธิบดีกำหนด"}
\end{quote}

หักลดหย่อนได้ตามจำนวนที่จ่ายจริง แต่ไม่เกิน \hilight{15,000 บาท}

\subsection*{สรุปเงื่อนไขต่าง ๆ ที่สำคัญกรณีค่าลดหย่อนเบี้ยประกันสุขภาพบิดามารดา}

\begin{enumerate}
    \item \textbf{ค่าลดหย่อน:} ให้หักลดหย่อนเบี้ยประกันสุขภาพสำหรับ "บิดาของผู้มีเงินได้" "มารดาของผู้มีเงินได้" "บิดาของคู่สมรสของผู้มีเงินได้" และ "มารดาของคู่สมรสของผู้มีเงินได้" ตามที่จ่ายจริงรวมแล้วไม่เกิน 15,000 บาท
    \item \textbf{อายุ:} ไม่มีเงื่อนไขเรื่องอายุของบิดามารดาฯ (แตกต่างจากค่าลดหย่อนอุปการะเลี้ยงดูที่ต้องอายุ 60 ปีขึ้นไป)
    \item \textbf{ฐานะ:} เงินได้พึงประเมินซึ่งไม่ได้รับการยกเว้นภาษีที่บิดามารดาฯ แต่ละคนได้รับในแต่ละปีภาษีไม่เกิน 30,000 บาท
    \item \textbf{บุตรชอบด้วยกฎหมาย:} ผู้มีเงินได้หรือสามีหรือภริยาของผู้มีเงินได้จะต้องเป็นบุตรชอบด้วยกฎหมายของบิดามารดาที่ผู้มีเงินได้ใช้สิทธิหักลดหย่อน (บุตรบุญธรรมไม่มีสิทธิ)
    \item \textbf{การแจ้งความประสงค์:} แจ้งความประสงค์ที่จะใช้สิทธิยกเว้นภาษีเงินได้ต่อบริษัทประกันชีวิตหรือบริษัทประกันวินาศภัยที่ได้เอาประกันไว้
    \item \textbf{กรณีผู้มีเงินได้หลายคนร่วมกันทำประกันสุขภาพสำหรับบิดามารดาฯ:} ให้ผู้มีเงินได้ทุกคนได้รับยกเว้นภาษีเงินได้โดยเฉลี่ยเบี้ยประกันภัยที่ผู้มีเงินได้ร่วมกันจ่ายจริงแต่ไม่เกินหนึ่งหมื่นห้าพันบาท ตามส่วนจำนวนผู้มีเงินได้
    \item \textbf{กรณีผู้มีเงินได้มิได้เป็นผู้อยู่ในประเทศไทย:} ให้หักลดหย่อนเบี้ยประกันสุขภาพสำหรับบิดามารดาได้เฉพาะการทำประกันภัยสำหรับบิดามารดาที่อยู่ในประเทศไทย
\end{enumerate}

\newpage
\section{ค่าลดหย่อนการอุปการะเลี้ยงดูคนพิการหรือคนทุพพลภาพ}

ฐานทางกฎหมาย: มาตรา 47(1)(ฎ) แห่งประมวลรัษฎากร

หักลดหย่อนให้คนละ \hilight{60,000 บาท}

มาตรา 47(1)(ฎ) บัญญัติว่า

\begin{quote}
{\itshape "(ฎ) ค่าอุปการะเลี้ยงดูบิดามารดา สามีหรือภริยา บุตรชอบด้วยกฎหมายหรือบุตรบุญธรรมของผู้มีเงินได้ บิดามารดาหรือบุตรชอบด้วยกฎหมายของสามีหรือภริยาของผู้มีเงินได้ หรือบุคคลอื่นที่ผู้มีเงินได้เป็นผู้ดูแลตามกฎหมายว่าด้วยการส่งเสริมและพัฒนาคุณภาพชีวิตคนพิการ คนละหกหมื่นบาท โดยบุคคลดังกล่าวต้องเป็นคนพิการซึ่งมีบัตรประจำตัวคนพิการตามกฎหมายว่าด้วยการส่งเสริมและพัฒนาคุณภาพชีวิตคนพิการ หรือเป็นคนทุพพลภาพ มีรายได้ไม่เพียงพอแก่การยังชีพ และอยู่ในความอุปการะเลี้ยงดูของผู้มีเงินได้ ทั้งนี้ตามหลักเกณฑ์ วิธีการ และเงื่อนไข รวมทั้งจำนวนคนพิการและคนทุพพลภาพในความอุปการะเลี้ยงดูของผู้มีเงินได้ที่อธิบดีประกาศกำหนด

การหักลดหย่อนสำหรับบุตรบุญธรรม ให้หักได้ในฐานะบุตรบุญธรรมเพียงฐานะเดียว

ค่าอุปการะเลี้ยงดูตามวรรคหนึ่ง ให้หักลดหย่อนสำหรับเงินได้พึงประเมินประจำปี พ.ศ. 2552 ที่จะต้องยื่นรายการในปี พ.ศ. 2553 เป็นต้นไป"} (พระราชบัญญัติแก้ไขเพิ่มเติมประมวลรัษฎากร ฉบับที่ 37 พ.ศ. 2552 ใช้บังคับตั้งแต่วันถัดจากวันประกาศในราชกิจจานุเบกษาเป็นต้นไป) (ดูประกาศอธิบดีกรมสรรพากรเกี่ยวกับภาษีเงินได้ ฉบับที่ 182)
\end{quote}

\subsection*{สรุปเงื่อนไขต่าง ๆ ที่สำคัญกรณีค่าอุปการะเลี้ยงดูคนพิการและคนทุพพลภาพ}

\begin{enumerate}
    \item \textbf{ค่าลดหย่อน:} ให้หักลดหย่อนค่าอุปการะเลี้ยงดูคนพิการ/คนทุพพลภาพที่เป็น "บิดามารดาของผู้มีเงินได้" "บิดามารดาของคู่สมรสของผู้มีเงินได้" "คู่สมรสของผู้มีเงินได้" "บุตรชอบด้วยกฎหมายหรือบุตรบุญธรรมของผู้มีเงินได้" "บุตรชอบด้วยกฎหมายของคู่สมรส" หรือ "บุคคลอื่นที่ผู้มีเงินได้เป็นผู้ดูแลตามกฎหมายว่าด้วยการส่งเสริมและพัฒนาคุณภาพชีวิตคนพิการจำนวน 1 คน" คนละ 60,000 บาท
    \item \textbf{อายุ:} ไม่มีเงื่อนไขเรื่องอายุ
    \item \textbf{ฐานะ:} เงินได้พึงประเมินซึ่งไม่ได้รับการยกเว้นภาษีที่คนพิการหรือคนทุพพลภาพได้รับในแต่ละปีภาษีไม่เกิน 30,000 บาท
    \item \textbf{บุตรชอบด้วยกฎหมายและบุตรบุญธรรม:} หักในฐานะบุตรบุญธรรมได้ฐานะเดียว
    \item \textbf{คนพิการและคนทุพพลภาพ:} คนพิการที่ดูแลต้องมีบัตรประจำตัวคนพิการตามกฎหมายฯ และผู้มีเงินได้เป็นผู้ดูแลที่มีชื่อในบัตรประจำตัว ส่วนคนทุพพลภาพต้องมีใบรับรองแพทย์ที่ขึ้นทะเบียนรับใบอนุญาตรับรองว่ามีภาวะจำกัดหรือขาดความสามารถในการประกอบกิจวัตรหลักอันเป็นปกติเยี่ยงบุคคลทั่วไปอันเนื่องมาจากสาเหตุทางปัญหาสุขภาพหรือความเจ็บป่วยที่เป็นต่อเนื่องมาไม่น้อยกว่า 180 วัน หรือทุพพลภาพมาแล้วไม่น้อยกว่า 180 วัน และผู้มีเงินได้มีหนังสือรับรองเป็นผู้อุปการะเลี้ยงดูคนทุพพลภาพ
    \item \textbf{กรณีเป็นทั้งคนพิการและคนทุพพลภาพ:} ให้ผู้มีเงินได้หักในฐานะคนพิการได้เพียงฐานะเดียว
    \item \textbf{กรณีผู้มีเงินได้มิได้เป็นผู้อยู่ในประเทศไทย:} ให้หักลดหย่อนได้เฉพาะคนพิการและคนทุพพลภาพที่อยู่ในประเทศไทย
\end{enumerate}

\newpage
\section{ค่าลดหย่อนกรณีผู้ถึงแก่ความตายระหว่างปีภาษี}

ฐานทางกฎหมาย: มาตรา 47(4) แห่งประมวลรัษฎากร

ในกรณีผู้มีเงินได้ถึงแก่ความตายระหว่างปีภาษี ให้หักลดหย่อนได้ \hilight{เสมือนว่ามีชีวิตอยู่ตลอดปีภาษี} ที่ผู้นั้นถึงแก่ความตาย กล่าวคือ ไม่มีการลดสัดส่วนค่าลดหย่อนตามระยะเวลาที่มีชีวิตอยู่จริงในปีภาษีนั้น แต่ให้หักลดหย่อนได้เต็มจำนวนทุกรายการที่ผู้ถึงแก่ความตายมีสิทธิได้รับ เสมือนมีชีวิตอยู่จนสิ้นปีภาษี

\section{ค่าลดหย่อนกรณีกองมรดกที่ยังไม่ได้แบ่ง}

ฐานทางกฎหมาย: มาตรา 47(5) แห่งประมวลรัษฎากร

กองมรดกที่ยังไม่ได้แบ่งซึ่งมีหน้าที่ต้องเสียภาษีเงินได้บุคคลธรรมดาในฐานะหน่วยภาษีตามกฎหมาย ให้หักลดหย่อนได้ \hilight{60,000 บาท} ซึ่งเทียบเท่ากับค่าลดหย่อนส่วนตัวของบุคคลธรรมดาทั่วไป

\section{ค่าลดหย่อนกรณีห้างหุ้นส่วนสามัญหรือคณะบุคคลที่มิใช่นิติบุคคล}

ฐานทางกฎหมาย: มาตรา 47(6) แห่งประมวลรัษฎากร

ให้หักลดหย่อนได้ตามมาตรา 47(1)(ก) สำหรับผู้เป็นหุ้นส่วนหรือบุคคลในคณะบุคคลแต่ละคนซึ่งเป็นผู้อยู่ในประเทศไทย แต่รวมกันต้องไม่เกิน \hilight{120,000 บาท}

กล่าวคือ แม้ห้างหุ้นส่วนสามัญหรือคณะบุคคลที่มิใช่นิติบุคคลจะมีผู้เป็นหุ้นส่วนหรือบุคคลในคณะมากกว่า 2 คน แต่เพดานรวมของค่าลดหย่อนตามมาตรา 47(1)(ก) สำหรับหน่วยภาษีนี้ถูกจำกัดไว้ที่ 120,000 บาทเสมอ (เทียบเท่าค่าลดหย่อนส่วนตัวของสามีภริยาตามมาตรา 47(2))

\newpage
\section{ค่าลดหย่อนเงินบริจาค}

ฐานทางกฎหมาย: มาตรา 47(1)(ฏ) และมาตรา 47(7) แห่งประมวลรัษฎากร

\subsection*{เงินบริจาคแก่พรรคการเมือง (มาตรา 47(1)(ฏ))}

\begin{quote}
\textit{"(ฏ) เงินที่บริจาคแก่พรรคการเมือง หรือเงิน ทรัพย์สิน หรือประโยชน์อื่นใดที่ให้เพื่อสนับสนุนการจัดกิจกรรมระดมทุนของพรรคการเมืองตามกฎหมายประกอบรัฐธรรมนูญว่าด้วยพรรคการเมือง ตามจำนวนที่จ่ายจริงแต่รวมกันไม่เกินหนึ่งหมื่นบาท ทั้งนี้ตามหลักเกณฑ์ วิธีการ และเงื่อนไขที่อธิบดีประกาศกำหนด"} (แก้ไขเพิ่มเติมโดยพระราชบัญญัติแก้ไขเพิ่มเติมประมวลรัษฎากร ฉบับที่ 51 พ.ศ. 2562 ให้ใช้บังคับกับการคำนวณภาษีสำหรับเงินได้พึงประเมินประจำปี พ.ศ. 2561 ที่จะต้องยื่นรายการในปี พ.ศ. 2562 เป็นต้นไป)
\end{quote}

หักลดหย่อนได้ตามที่จ่ายจริง แต่ไม่เกิน \hilight{10,000 บาท}

\subsection*{เงินบริจาคทั่วไป (มาตรา 47(7))}

\begin{quote}
{\itshape "(7) เมื่อได้หักลดหย่อนตาม (1)(2)(3)(4)(5) หรือ (6) แล้ว เหลือเท่าใดให้หักลดหย่อนได้อีกสำหรับเงินบริจาคดังต่อไปนี้ โดยให้หักได้เท่าจำนวนที่บริจาค แต่ต้องไม่เกินร้อยละ 10 ของเงินที่เหลือนั้น

(ก) เงินที่บริจาคแก่สถานพยาบาลและสถานศึกษาของทางราชการ

(ข) เงินที่บริจาคเป็นสาธารณประโยชน์แก่องค์การหรือสถานสาธารณกุศล หรือแก่สถานพยาบาลและสถานศึกษาอื่นนอกจากที่กล่าวใน (ก) ทั้งนี้ตามที่รัฐมนตรีประกาศกำหนดในราชกิจจานุเบกษา"} (พระราชบัญญัติแก้ไขเพิ่มเติม ฉบับที่ 19 พ.ศ. 2508 ใช้บังคับปีภาษี 2509 เป็นต้นไป) (ดูพระราชกฤษฎีกาฉบับที่ 420 พ.ศ. 2537) (ดูพระราชกฤษฎีกาฉบับที่ 274 พ.ศ. 2537) (ดูพระราชกฤษฎีกาฉบับที่ 317 พ.ศ. 2541) (ดูประกาศกระทรวงการคลัง ว่าด้วยภาษีเงินได้และภาษีมูลค่าเพิ่ม ฉบับที่ 2) (ดูคำสั่งกรมสรรพากรที่ ป.120/2545)
\end{quote}

\subsection*{กรณียกเว้นเงินได้หนึ่งเท่าหรือสองเท่าของจำนวนเงินที่บริจาค}

นอกจากเงินบริจาคทั่วไปตามมาตรา 47(7) ที่หักลดหย่อนได้เท่าจำนวนที่บริจาคจริงแต่ไม่เกินร้อยละ 10 ของเงินได้พึงประเมินหลังหักค่าใช้จ่ายและค่าลดหย่อนแล้ว ยังมีพระราชกฤษฎีกาอีกหลายฉบับที่ให้สิทธิ \hilight{ยกเว้นภาษีเป็น 1 เท่าหรือ 2 เท่า} ของจำนวนเงินที่บริจาคจริง สำหรับการบริจาคเพื่อวัตถุประสงค์เฉพาะที่รัฐต้องการส่งเสริม ซึ่งถือเป็นรูปแบบ Exemption as Deduction ดังที่กล่าวไว้ในข้อสังเกตตอนต้นของเอกสารนี้

\textbf{พระราชกฤษฎีกาฉบับที่ 619} ออกตามความมาตรา 3(1) ให้ยกเว้นภาษีเงินได้สำหรับการบริจาคให้แก่กองทุนพัฒนาเทคโนโลยีเพื่อการศึกษาที่กระทรวงศึกษาธิการจัดตั้งขึ้น สำหรับบุคคลธรรมดา ให้ยกเว้นสำหรับเงินได้พึงประเมินหลังจากหักค่าใช้จ่ายและหักลดหย่อนตามมาตรา 47(1)(2)(3)(4)(5) หรือ (6) แห่งประมวลรัษฎากร เท่าจำนวนเงินที่บริจาค แต่เมื่อรวมกับเงินบริจาคตามมาตรา 47(7) แห่งประมวลรัษฎากรแล้ว ต้องไม่เกินร้อยละ 10 ของเงินได้พึงประเมินหลังจากหักค่าใช้จ่ายและหักลดหย่อนนั้น

\textbf{พระราชกฤษฎีกาฉบับที่ 519} ออกตามความมาตรา 3(1) ให้ยกเว้นภาษีเงินได้สำหรับเงินได้ที่จ่ายเป็นค่าใช้จ่ายในการจัดให้คนพิการได้รับสิทธิเข้าถึงและใช้ประโยชน์ได้จากสิ่งอำนวยความสะดวกอันเป็นสาธารณะ ตลอดจนสวัสดิการและความช่วยเหลืออื่นจากรัฐตามกฎหมายว่าด้วยการส่งเสริมและพัฒนาคุณภาพชีวิตคนพิการ สำหรับบุคคลธรรมดา ให้ยกเว้นภาษีสำหรับเงินได้พึงประเมินหลังจากหักค่าใช้จ่ายและหักลดหย่อนตามมาตรา 47(1)(2)(3)(4)(5) หรือ (6) แห่งประมวลรัษฎากร เป็นจำนวนร้อยละหนึ่งร้อยของเงินได้ที่จ่ายเป็นค่าใช้จ่ายในการจัดให้คนพิการได้รับสิทธิประโยชน์ แต่เมื่อรวมกับเงินได้ที่ได้รับยกเว้นสำหรับการจ่ายเป็นค่าใช้จ่ายเพื่อสนับสนุนการศึกษาสำหรับโครงการที่กระทรวงศึกษาธิการให้ความเห็นชอบแล้ว ต้องไม่เกินร้อยละ 10 ของเงินได้พึงประเมินหลังจากหักค่าใช้จ่ายและหักลดหย่อนดังกล่าว

\textbf{พระราชกฤษฎีกาฉบับที่ 420} ออกตามความมาตรา 3(1) ให้ยกเว้นภาษีเงินได้สำหรับเงินได้ที่จ่ายเป็นค่าใช้จ่ายเพื่อสนับสนุนการศึกษาให้แก่สถานศึกษาของทางราชการ สถานศึกษาขององค์การของรัฐบาล โรงเรียนเอกชนที่ตั้งขึ้นตามกฎหมายว่าด้วยโรงเรียนเอกชน สถาบันอุดมศึกษาเอกชนที่ตั้งขึ้นตามกฎหมายว่าด้วยสถาบันอุดมศึกษาเอกชน หรือสถาบันอุดมศึกษาซึ่งคณะกรรมการพัฒนาการจัดการศึกษาโดยสถาบันอุดมศึกษาที่มีศักยภาพสูงจากต่างประเทศอนุมัติ สำหรับบุคคลธรรมดา ให้ยกเว้นภาษีเงินได้สำหรับเงินได้พึงประเมินหลังจากหักค่าใช้จ่ายและหักค่าลดหย่อนตามมาตรา 47(1)(2)(3)(4)(5) หรือ (6) แห่งประมวลรัษฎากร เป็นจำนวน \hilight{สองเท่า} ของรายจ่ายที่จ่ายไปเป็นค่าใช้จ่ายเพื่อสนับสนุนการศึกษา แต่ต้องไม่เกินร้อยละ 10 ของเงินได้พึงประเมินหลังจากหักค่าใช้จ่ายและหักค่าลดหย่อนดังกล่าวนั้น ค่าใช้จ่ายเพื่อสนับสนุนการศึกษาต้องเป็นค่าใช้จ่ายสำหรับโครงการที่กระทรวงศึกษาธิการให้ความเห็นชอบ และเป็นค่าใช้จ่ายสำหรับรายการดังต่อไปนี้ (1) จัดหาหรือจัดสร้างอาคาร อาคารพร้อมที่ดิน หรือที่ดินให้แก่สถานศึกษาเพื่อใช้ประโยชน์ทางการศึกษา (2) จัดหาวัสดุอุปกรณ์เพื่อการศึกษา แบบเรียน ตำรา หนังสือทางวิชาการ สื่อและเทคโนโลยีเพื่อการศึกษา ตลอดจนวัสดุอุปกรณ์อื่น ๆ ที่เกี่ยวข้องกับการศึกษาให้แก่สถานศึกษาตามที่รัฐมนตรีว่าการกระทรวงการคลังกำหนด หรือ (3) จัดหาครู อาจารย์ หรือผู้ทรงคุณวุฒิทางการศึกษา หรือเป็นทุนการศึกษา การประดิษฐ์ การพัฒนา การค้นคว้า หรือการวิจัย สำหรับนักเรียน นิสิต หรือนักศึกษาของสถานศึกษา

\textbf{พระราชกฤษฎีกาฉบับที่ 651} ออกตามความมาตรา 3(1) ให้ยกเว้นภาษีเงินได้สำหรับการบริจาคให้แก่กองทุนยุติธรรมตามกฎหมายว่าด้วยกองทุนยุติธรรม สำหรับบุคคลธรรมดา ให้ยกเว้นสำหรับเงินได้พึงประเมินหลังจากหักค่าใช้จ่ายและหักลดหย่อนตามมาตรา 47(1)(2)(3)(4)(5) หรือ (6) แห่งประมวลรัษฎากร เป็นจำนวน \hilight{สองเท่า} ของจำนวนเงินที่บริจาค แต่เมื่อรวมกับเงินได้ที่ได้รับยกเว้นสำหรับการจ่ายเป็นค่าใช้จ่ายเพื่อสนับสนุนการศึกษาสำหรับโครงการที่กระทรวงศึกษาธิการให้ความเห็นชอบแล้ว ต้องไม่เกินร้อยละ 10 ของเงินได้พึงประเมินหลังจากหักค่าใช้จ่ายและหักลดหย่อนนั้น

\textbf{พระราชกฤษฎีกาฉบับที่ 551} ออกตามความมาตรา 3(1) ให้ยกเว้นภาษีเงินได้สำหรับเงินได้ที่บุคคลธรรมดาได้จ่ายให้แก่สถานศึกษาของทางราชการ สถานศึกษาขององค์การของรัฐบาล โรงเรียนเอกชนที่ตั้งขึ้นตามกฎหมายว่าด้วยโรงเรียนเอกชน หรือสถาบันอุดมศึกษาเอกชนที่ตั้งขึ้นตามกฎหมายว่าด้วยสถาบันอุดมศึกษาเอกชน เพื่อใช้ในการจัดหาหนังสือหรือสื่ออิเล็กทรอนิกส์เพื่อส่งเสริมการอ่าน สำหรับบุคคลธรรมดา ให้ยกเว้นสำหรับเงินได้พึงประเมินหลังจากหักค่าใช้จ่ายและหักลดหย่อนตามมาตรา 47(1)(2)(3)(4)(5) หรือ (6) แห่งประมวลรัษฎากร เป็นจำนวน \hilight{สองเท่า} ของรายจ่ายที่จ่ายไป แต่เมื่อรวมกับรายจ่ายที่จ่ายไปเป็นค่าใช้จ่ายเพื่อสนับสนุนการศึกษาสำหรับโครงการที่กระทรวงศึกษาธิการให้ความเห็นชอบแล้ว ต้องไม่เกินร้อยละ 10 ของเงินได้พึงประเมินหลังจากหักค่าใช้จ่ายและหักลดหย่อนนั้น

\textbf{พระราชกฤษฎีกาฉบับที่ 735} ออกตามความมาตรา 6 ให้ยกเว้นภาษีเงินได้สำหรับเงินได้ที่ได้จ่ายเพื่อการลงทุนในหุ้นหรือการเป็นหุ้นส่วนเพื่อการจัดตั้งหรือการเพิ่มทุนของวิสาหกิจเพื่อสังคมที่ได้จดแจ้งต่ออธิบดีตามมาตรา 11 แล้ว สำหรับบุคคลธรรมดา ให้ยกเว้นสำหรับเงินได้พึงประเมินเท่าที่ได้จ่ายไปเพื่อการลงทุนในหุ้นหรือการเป็นหุ้นส่วนเพื่อการจัดตั้งหรือเพื่อการเพิ่มทุน แล้วแต่กรณี แต่เมื่อรวมกันแล้วต้องไม่เกินกรณีละหนึ่งแสนบาทสำหรับปีภาษีนั้น และตามมาตรา 8 ให้ยกเว้นภาษีเงินได้สำหรับการบริจาคผ่านระบบบริจาคอิเล็กทรอนิกส์ให้แก่กองทุนส่งเสริมวิสาหกิจเพื่อสังคมที่ได้กระทำตั้งแต่วันที่พระราชกฤษฎีกานี้มีผลใช้บังคับถึงวันที่ 31 ธันวาคม พ.ศ. 2566 สำหรับบุคคลธรรมดา ให้ยกเว้นสำหรับเงินได้พึงประเมินหลังจากหักค่าใช้จ่ายและหักลดหย่อนตามมาตรา 47(1)(2)(3)(4)(5) หรือ (6) แห่งประมวลรัษฎากร เท่าจำนวนเงินที่บริจาค แต่เมื่อรวมกับเงินบริจาคตามมาตรา 47(7) แห่งประมวลรัษฎากรแล้ว ต้องไม่เกินร้อยละ 10 ของเงินได้พึงประเมินหลังหักค่าใช้จ่ายและหักลดหย่อนนั้น

\subsection*{เงินบริจาคเฉพาะประเภทตามกฎกระทรวงฉบับที่ 126 ที่หักได้เท่าจำนวนที่บริจาค}

นอกจากนี้ ยังมีเงินบริจาคบางประเภทที่กฎกระทรวงฉบับที่ 126 ให้สิทธิยกเว้นภาษีเงินได้เท่าจำนวนเงินที่บริจาค (ไม่ใช่ 2 เท่า) โดยเมื่อรวมกับเงินบริจาคตามมาตรา 47(7) แล้วต้องไม่เกินร้อยละ 10 ของเงินได้พึงประเมินหลังจากหักค่าใช้จ่ายและหักลดหย่อนดังกล่าวด้วยเช่นกัน ได้แก่

\textbf{ข้อ 2(68)} เงินได้พึงประเมินหลังจากหักค่าใช้จ่ายและหักลดหย่อนตามมาตรา 47(1)(2)(3)(4)(5) หรือ (6) แห่งประมวลรัษฎากร เท่าจำนวนเงินที่ได้บริจาคให้แก่การกีฬาแห่งประเทศไทย เพื่อส่งเสริมการกีฬา คณะกรรมการกีฬาจังหวัดที่จัดตั้งขึ้นตามกฎหมายว่าด้วยการกีฬาแห่งประเทศไทยเพื่อส่งเสริมกีฬาในจังหวัด กรมพลศึกษาเพื่อการจัดการแข่งขันกีฬานักเรียน หรือสมาคมกีฬาจังหวัดหรือสมาคมกีฬาแห่งประเทศไทยที่จัดตั้งขึ้นโดยได้รับอนุญาตจากการกีฬาแห่งประเทศไทยเพื่อการกีฬา แต่เมื่อรวมกับเงินบริจาคตามมาตรา 47(7) แห่งประมวลรัษฎากรแล้ว ต้องไม่เกินร้อยละ 10 ของเงินได้พึงประเมินหลังจากหักค่าใช้จ่ายและหักลดหย่อนดังกล่าว (แก้ไขเพิ่มเติมโดยกฎกระทรวงฉบับที่ 294 พ.ศ. 2555 ใช้บังคับ 26 ธันวาคม พ.ศ. 2555 เป็นต้นไป)

\textbf{ข้อ 2(70)} เงินได้พึงประเมินหลังจากหักค่าใช้จ่ายและค่าลดหย่อนตามมาตรา 47(1)(2)(3)(4)(5) หรือ (6) แห่งประมวลรัษฎากร เท่าจำนวนที่บริจาคให้แก่ส่วนราชการ เพื่อช่วยเหลือผู้ประสบอุทกภัย วาตภัย อัคคีภัย หรือภัยธรรมชาติอื่น แต่เมื่อรวมกับเงินบริจาคตามมาตรา 47(7) แห่งประมวลรัษฎากรแล้ว ต้องไม่เกินร้อยละ 10 ของเงินได้พึงประเมินหลังจากหักค่าใช้จ่ายและหักค่าลดหย่อนดังกล่าวนั้น ทั้งนี้สำหรับเงินได้พึงประเมินประจำปี พ.ศ. 2547 ที่ต้องยื่นรายการในปี พ.ศ. 2548 เป็นต้นไป (แก้ไขเพิ่มเติมโดยกฎกระทรวงฉบับที่ 253 พ.ศ. 2548)

\textbf{ข้อ 2(88)} เงินได้พึงประเมินหลังจากหักค่าใช้จ่ายและหักลดหย่อนตามมาตรา 47(1)(2)(3)(4)(5) หรือ (6) แห่งประมวลรัษฎากร เท่าจำนวนเงินที่บริจาคให้แก่กรมศิลปากร เพื่อการบูรณะโบราณสถาน โบราณวัตถุ และศิลปวัตถุตามกฎหมายว่าด้วยโบราณสถาน โบราณวัตถุ ศิลปวัตถุ และพิพิธภัณฑสถานแห่งชาติ แต่เมื่อรวมกับเงินบริจาคตามมาตรา 47(7) แห่งประมวลรัษฎากรแล้วต้องไม่เกินร้อยละ 10 ของเงินได้พึงประเมินหลังจากหักค่าใช้จ่ายและหักลดหย่อนดังกล่าว ทั้งนี้สำหรับเงินได้พึงประเมินประจำปี พ.ศ. 2556 ที่ต้องยื่นรายการในปี พ.ศ. 2557 เป็นต้นไป (แก้ไขเพิ่มเติมโดยกฎกระทรวงฉบับที่ 302 พ.ศ. 2556 ใช้บังคับ 16 ธันวาคม พ.ศ. 2556 เป็นต้นไป)

\hilight{ข้อสังเกตสำคัญสำหรับการทำข้อสอบ:} ต้องแยกแยะระหว่าง (1) เงินบริจาคที่หักลดหย่อนตามมาตรา 47(7) แบบทั่วไป ซึ่งมีเพดานรวมไม่เกินร้อยละ 10 ของเงินได้หลังหักค่าใช้จ่ายและค่าลดหย่อน กับ (2) เงินบริจาคที่ได้รับยกเว้นภาษีเป็น 1 เท่าหรือ 2 เท่าของจำนวนที่บริจาคตามพระราชกฤษฎีกาเฉพาะฉบับ ซึ่งเป็นรูปแบบ "ยกเว้นเงินได้เท่าที่ได้จ่าย" (Exemption as Deduction) มิใช่การหักลดหย่อนตามมาตรา 47 โดยตรง แต่ทั้งสองกรณีมักมีเพดานรวมกันไม่เกินร้อยละ 10 ของเงินได้พึงประเมินหลังหักค่าใช้จ่ายและค่าลดหย่อนเช่นเดียวกัน เมื่อนำเงินบริจาคทุกประเภทมารวมกัน

\newpage
\section{ตัวอย่างค่าลดหย่อนและข้อยกเว้นภาษีอื่น ๆ ที่ไม่ผ่านมาตรา 47 โดยตรง}

นอกจากค่าลดหย่อนตามมาตรา 47 และเงินบริจาคที่กล่าวมาแล้ว ยังมีพระราชกฤษฎีกาและกฎกระทรวงอีกหลายฉบับที่ใช้รูปแบบ "ยกเว้นเงินได้เท่าที่ได้จ่าย" (Exemption as Deduction) สำหรับนโยบายเฉพาะด้านของรัฐ ซึ่งมีลักษณะคล้ายค่าลดหย่อนในทางปฏิบัติ ตัวอย่างที่สำคัญมีดังนี้

\subsection*{พระราชกฤษฎีกาฉบับที่ 728: เขตพัฒนาพิเศษเฉพาะกิจ}

ออกตามความมาตรา 3(1) นิยาม "เขตพัฒนาพิเศษเฉพาะกิจ" หมายความว่า ท้องที่จังหวัดนราธิวาส จังหวัดปัตตานี จังหวัดยะลา จังหวัดสงขลาเฉพาะในท้องที่อำเภอจะนะ อำเภอเทพา อำเภอนาทวี และอำเภอสะบ้าย้อย และจังหวัดสตูล

มาตรา 4 ให้ยกเว้นภาษีเงินได้ตามส่วน 2 และส่วน 3 หมวด 3 ในลักษณะ 2 แห่งประมวลรัษฎากร เป็นจำนวนร้อยละหนึ่งร้อยของเงินได้เท่าที่ได้จ่ายเป็นค่าซื้อและค่าติดตั้งระบบกล้องโทรทัศน์วงจรปิด ณ สถานประกอบกิจการที่ตั้งอยู่ในเขตพัฒนาพิเศษเฉพาะกิจ ดังต่อไปนี้ (1) สำหรับบุคคลธรรมดาซึ่งมีเงินได้พึงประเมินตามมาตรา 40(5)(6)(7) และ (8) แห่งประมวลรัษฎากร และมีสถานประกอบกิจการตั้งอยู่ในเขตพัฒนาพิเศษเฉพาะกิจ ที่ได้จ่ายค่าซื้อและค่าติดตั้งระบบกล้องโทรทัศน์วงจรปิดดังกล่าวตั้งแต่วันที่ 1 มกราคม พ.ศ. 2564 ถึงวันที่ 31 ธันวาคม พ.ศ. 2566

\subsection*{พระราชกฤษฎีกาฉบับที่ 805: ระบบผลิตไฟฟ้าจากพลังงานแสงอาทิตย์}

ออกตามความมาตรา 3(1) มาตรา 3 ให้ยกเว้นภาษีเงินได้ตามส่วน 2 หมวด 3 ในลักษณะ 2 แห่งประมวลรัษฎากร ให้แก่บุคคลธรรมดาแต่ไม่รวมถึงห้างหุ้นส่วนสามัญ คณะบุคคลที่ไม่ใช่นิติบุคคล หรือกองมรดกที่ยังมิได้แบ่ง สำหรับเงินได้เท่าที่ได้จ่ายเป็นค่าซื้ออุปกรณ์และค่าติดตั้งระบบการผลิตไฟฟ้าจากพลังงานแสงอาทิตย์ที่ติดตั้งบนหลังคา ดาดฟ้า หรือส่วนหนึ่งส่วนใดบนอาคารซึ่งบุคคลอาจเข้าอยู่หรือใช้สอยได้ ซึ่งเชื่อมต่อกับระบบโครงข่ายไฟฟ้าของการไฟฟ้านครหลวงหรือการไฟฟ้าส่วนภูมิภาค ตามจำนวนที่จ่ายจริงแต่รวมกันแล้วไม่เกิน \hilight{สองแสนบาท} การได้รับสิทธิยกเว้นภาษีเงินได้ตามวรรคหนึ่งต้องเป็นค่าซื้ออุปกรณ์และค่าติดตั้งระบบการผลิตไฟฟ้าจากพลังงานแสงอาทิตย์จำนวนหนึ่งระบบที่ได้จ่ายไปตั้งแต่วันที่พระราชกฤษฎีกานี้ใช้บังคับถึงวันที่ 31 ธันวาคม พ.ศ. 2571 ทั้งนี้สามารถใช้สิทธิยกเว้นภาษีเงินได้ได้หนึ่งครั้งในปีภาษีที่เชื่อมต่อกับระบบโครงข่ายไฟฟ้าสำเร็จ รวมทั้งต้องไม่นำค่าซื้ออุปกรณ์และค่าติดตั้งระบบดังกล่าวไปใช้สิทธิยกเว้นภาษีเงินได้ตามมาตรา 4 ไม่ว่าทั้งหมดหรือบางส่วน

\subsection*{กฎกระทรวงเกี่ยวกับมาตรการกระตุ้นเศรษฐกิจประเภท "ช้อปดีมีคืน"}

กฎกระทรวงชุดนี้ออกตามความมาตรา 42(17) แห่งประมวลรัษฎากร เป็นมาตรการชั่วคราวที่ออกต่อเนื่องเป็นรายปีเพื่อกระตุ้นการใช้จ่ายของประชาชน โดยมีรายละเอียดแต่ละฉบับดังนี้

\textbf{กฎกระทรวงฉบับที่ 344} ข้อ 2 กำหนดให้เงินได้ดังต่อไปนี้เป็นเงินได้ที่ได้รับยกเว้นไม่ต้องรวมคำนวณเพื่อเสียภาษีเงินได้บุคคลธรรมดา (1) เงินได้เท่าที่ผู้มีเงินได้ได้จ่ายเป็นค่าบริการให้แก่ผู้ประกอบธุรกิจนำเที่ยวตามกฎหมายว่าด้วยธุรกิจนำเที่ยวและมัคคุเทศก์ สำหรับการเดินทางท่องเที่ยวในจังหวัดท่องเที่ยวรอง ค่าที่พักโรงแรมตามกฎหมายว่าด้วยโรงแรม ค่าที่พักโฮมสเตย์ไทย หรือค่าที่พักในสถานที่พักที่ไม่เป็นโรงแรม ในจังหวัดท่องเที่ยวรอง ตามจำนวนที่จ่ายจริงแต่ไม่เกินสองหมื่นบาท หรือ (2) เงินได้เท่าที่ผู้มีเงินได้ได้จ่ายเป็นค่าบริการให้แก่ผู้ประกอบธุรกิจนำเที่ยวตามกฎหมายว่าด้วยธุรกิจนำเที่ยวและมัคคุเทศก์ สำหรับการเดินทางท่องเที่ยวในจังหวัดอื่นซึ่งมิใช่จังหวัดท่องเที่ยวรอง ค่าที่พักโรงแรมตามกฎหมายว่าด้วยโรงแรม ค่าที่พักโฮมสเตย์ไทย หรือค่าที่พักในสถานที่พักที่ไม่เป็นโรงแรมในจังหวัดอื่นซึ่งมิใช่จังหวัดท่องเที่ยวรอง ตามจำนวนที่จ่ายจริงแต่ไม่เกินหนึ่งหมื่นห้าพันบาท ในกรณีที่ผู้มีเงินได้ได้จ่ายค่าบริการทั้งสองกรณีรวมกันแล้วไม่เกินสองหมื่นบาท ทั้งนี้เงินได้ที่ได้รับยกเว้นต้องเป็นค่าบริการหรือค่าที่พักที่ได้จ่ายตั้งแต่วันที่ 30 เมษายน พ.ศ. 2562 ถึงวันที่ 30 มิถุนายน พ.ศ. 2562

\textbf{กฎกระทรวงฉบับที่ 379} กำหนดให้เงินได้เท่าที่ได้จ่ายเป็นค่าซื้อสินค้าหรือค่าบริการให้แก่ผู้ประกอบการจดทะเบียนภาษีมูลค่าเพิ่มและได้รับใบกำกับภาษีตามมาตรา 86/4 แห่งประมวลรัษฎากร สำหรับการซื้อสินค้าหรือรับบริการในราชอาณาจักร ตั้งแต่วันที่ 1 มกราคม พ.ศ. 2565 ถึงวันที่ 15 กุมภาพันธ์ พ.ศ. 2565 ตามจำนวนที่จ่ายจริงแต่ไม่เกินสามหมื่นบาท เป็นเงินได้พึงประเมินที่ได้รับยกเว้นไม่ต้องรวมคำนวณเพื่อเสียภาษีเงินได้บุคคลธรรมดา ทั้งนี้ ตามหลักเกณฑ์ วิธีการ และเงื่อนไขที่อธิบดีประกาศกำหนด

ค่าซื้อสินค้าหรือค่าบริการที่ไม่รวมอยู่ในสิทธิประโยชน์นี้ ได้แก่ ค่าซื้อสุรา เบียร์ และไวน์ ค่าซื้อยาสูบ ค่าซื้อน้ำมันและก๊าซสำหรับเติมยานพาหนะ ค่าซื้อรถยนต์ รถจักรยานยนต์ และเรือ ค่าสาธารณูปโภค ค่าน้ำประปา ค่าไฟฟ้า ค่าบริการสัญญาณโทรศัพท์ และค่าบริการสัญญาณอินเทอร์เน็ต ค่าบริการที่มีข้อตกลงการให้บริการและผู้รับบริการสามารถใช้บริการดังกล่าวนอกเหนือจากระยะเวลาที่กำหนด และค่าเบี้ยประกันวินาศภัย ผู้มีเงินได้ที่ได้รับสิทธิยกเว้นภาษีเงินได้ตามกฎกระทรวงนี้ต้องไม่เป็นห้างหุ้นส่วนสามัญหรือคณะบุคคลที่มิใช่นิติบุคคล

\textbf{กฎกระทรวงฉบับที่ 386} กำหนดสิทธิประโยชน์ในลักษณะเดียวกัน สำหรับการซื้อสินค้าหรือรับบริการในราชอาณาจักรตั้งแต่วันที่ 1 มกราคม พ.ศ. 2566 ถึงวันที่ 15 กุมภาพันธ์ พ.ศ. 2566 ตามจำนวนที่จ่ายจริงแต่ไม่เกินสามหมื่นบาท โดยมีรายการที่ไม่รวมอยู่ในสิทธิประโยชน์เช่นเดียวกันกับฉบับที่ 379

\textbf{กฎกระทรวงฉบับที่ 391} กำหนดสิทธิประโยชน์สำหรับการซื้อสินค้าหรือรับบริการในราชอาณาจักรตั้งแต่วันที่ 1 มกราคม พ.ศ. 2567 ถึงวันที่ 15 กุมภาพันธ์ พ.ศ. 2567 ตามจำนวนที่จ่ายจริงแต่ไม่เกินห้าหมื่นบาท โดยกรณีที่ผู้มีเงินได้จ่ายค่าซื้อสินค้าหรือค่าบริการเกินสามหมื่นบาท ให้เงินได้เท่าที่ได้จ่ายเป็นค่าซื้อสินค้าหรือค่าบริการดังต่อไปนี้ ตามจำนวนที่จ่ายจริงในส่วนที่เกินสามหมื่นบาทแต่ไม่เกินสองหมื่นบาท เป็นเงินได้พึงประเมินที่ได้รับยกเว้นเพิ่มเติม ได้แก่ (1) ค่าซื้อสินค้าหนึ่งตำบลหนึ่งผลิตภัณฑ์ซึ่งเป็นสินค้าที่ได้ลงทะเบียนกับกรมการพัฒนาชุมชนแล้ว (2) ค่าซื้อสินค้าหรือค่าบริการที่จ่ายให้แก่วิสาหกิจชุมชนที่ได้จดทะเบียนตามกฎหมายว่าด้วยการส่งเสริมวิสาหกิจชุมชนต่อกรมส่งเสริมการเกษตร และ (3) ค่าซื้อสินค้าหรือค่าบริการที่จ่ายให้แก่วิสาหกิจเพื่อสังคมที่ได้จดทะเบียนตามกฎหมายว่าด้วยการส่งเสริมวิสาหกิจเพื่อสังคมต่อสำนักงานส่งเสริมวิสาหกิจเพื่อสังคม ทั้งนี้เฉพาะกรณีที่ผู้มีเงินได้ได้รับใบกำกับภาษีตามมาตรา 86/4 แห่งประมวลรัษฎากร หรือใบรับตามมาตรา 105 แห่งประมวลรัษฎากร ที่ได้จัดทำโดยวิธีการทางอิเล็กทรอนิกส์ตามมาตรา 3 โสฬส แห่งประมวลรัษฎากร

\textbf{กฎกระทรวงฉบับที่ 397} กำหนดสิทธิประโยชน์สำหรับการซื้อสินค้าหรือรับบริการในราชอาณาจักรตั้งแต่วันที่ 16 มกราคม พ.ศ. 2568 ถึงวันที่ 28 กุมภาพันธ์ พ.ศ. 2568 ตามจำนวนที่จ่ายจริงแต่ไม่เกินสามหมื่นบาท และในกรณีที่จ่ายเกินสามหมื่นบาท ให้เงินได้เท่าที่ได้จ่ายเป็นค่าซื้อสินค้าหนึ่งตำบลหนึ่งผลิตภัณฑ์ ค่าซื้อสินค้าหรือค่าบริการที่จ่ายให้แก่วิสาหกิจชุมชนที่ได้จดทะเบียนตามกฎหมายว่าด้วยการส่งเสริมวิสาหกิจชุมชนต่อกรมส่งเสริมการเกษตร หรือค่าซื้อสินค้าหรือค่าบริการที่จ่ายให้แก่วิสาหกิจเพื่อสังคมที่ได้จดทะเบียนตามกฎหมายว่าด้วยการส่งเสริมวิสาหกิจเพื่อสังคมต่อสำนักงานส่งเสริมวิสาหกิจเพื่อสังคม ตามจำนวนที่จ่ายจริงในส่วนที่เกินสามหมื่นบาทแต่ไม่เกินสองหมื่นบาท ได้รับยกเว้นเพิ่มเติมด้วย โดยฉบับนี้เพิ่มประเภทค่าที่พักโฮมสเตย์ที่ได้รับการรับรองมาตรฐานโฮมสเตย์ไทยจากกรมการท่องเที่ยว กระทรวงการท่องเที่ยวและกีฬา และค่าที่พักในสถานที่พักที่ไม่เป็นโรงแรมที่จ่ายให้แก่ผู้ประกอบกิจการสถานที่พักที่ไม่เป็นโรงแรมตามกฎหมายว่าด้วยโรงแรม เข้ามาในรายการที่ได้รับยกเว้นด้วย

\hilight{ข้อสังเกตสำคัญสำหรับมาตรการเหล่านี้:} ทุกฉบับกำหนดเงื่อนไขร่วมว่า (1) ต้องมีใบกำกับภาษีตามมาตรา 86/4 หรือใบรับตามมาตรา 105 ที่จัดทำโดยวิธีทางอิเล็กทรอนิกส์ผ่านระบบ e-Tax Invoice เท่านั้น (2) ผู้ได้รับสิทธิต้องไม่เป็นห้างหุ้นส่วนสามัญหรือคณะบุคคลที่มิใช่นิติบุคคล และ (3) มีรายการสินค้า/บริการที่ไม่เข้าเกณฑ์คล้ายกันในทุกฉบับ ได้แก่ สุรา เบียร์ ไวน์ ยาสูบ น้ำมันและแก๊สเติมยานพาหนะ รถยนต์ รถจักรยานยนต์ และเรือ ค่าสาธารณูปโภค (น้ำ ไฟ โทรศัพท์ อินเทอร์เน็ต) และค่าเบี้ยประกันวินาศภัย ข้อสอบมักถามเปรียบเทียบจำนวนเพดานและรายการที่ยกเว้น/ไม่ยกเว้นในแต่ละปี ผู้เรียนควรจดจำเพดานหลักของแต่ละปีและรายการข้อยกเว้นร่วมที่ปรากฏซ้ำทุกฉบับ

\newpage
\section{ตารางสรุปเพดานค่าลดหย่อนหลักทั้งหมด}

ตารางต่อไปนี้รวบรวมเพดานค่าลดหย่อนและข้อยกเว้นภาษีที่สำคัญที่สุดทั้งหมดที่กล่าวมาในเอกสารนี้ เพื่อใช้สำหรับท่องจำและทบทวนก่อนการทำข้อสอบ

\begin{longtable}{|p{6.5cm}|p{8cm}|}
\hline
\textbf{รายการค่าลดหย่อน} & \textbf{จำนวนสูงสุดที่หักได้} \\
\hline
\endhead
ค่าลดหย่อนส่วนตัว (มาตรา 47(1)(ก)) & 60,000 บาท \\
\hline
ค่าลดหย่อนคู่สมรส (มาตรา 47(1)(ข)) & 60,000 บาท (รวมกัน 120,000 บาท หากสามีภริยาต่างมีเงินได้) \\
\hline
ค่าลดหย่อนบุตรชอบด้วยกฎหมาย/บุตรบุญธรรม (มาตรา 47(1)(ค)) & คนละ 30,000 บาท (เพิ่มเป็นคนละ 60,000 บาท สำหรับบุตรชอบด้วยกฎหมายลำดับที่ 2 เป็นต้นไปที่เกิดตั้งแต่ปี 2561) \\
\hline
ค่าฝากครรภ์และคลอดบุตร (ข้อ 2(99)) & ไม่เกิน 60,000 บาท ต่อการตั้งครรภ์แต่ละคราว \\
\hline
เบี้ยประกันชีวิตของตนเอง (มาตรา 47(1)(ง) + ข้อ 2(61) วรรคหนึ่ง) & ไม่เกิน 100,000 บาท (10,000 + 90,000) \\
\hline
เบี้ยประกันชีวิตแบบบำนาญ (ข้อ 2(61) วรรคสอง) & เพิ่มอีกไม่เกิน 200,000 บาท (ไม่เกินร้อยละ 15 ของเงินได้พึงประเมิน) \\
\hline
เงินสะสมกองทุนสำรองเลี้ยงชีพ (มาตรา 47(1)(ช) + ข้อ 2(35)) & ไม่เกิน 500,000 บาท (10,000 + 490,000) \\
\hline
เงินฝากที่มีประกันชีวิต (ข้อ 2(94)) & ไม่เกิน 100,000 บาท \\
\hline
เบี้ยประกันสุขภาพตนเอง (ข้อ 2(97)) & ไม่เกิน 25,000 บาท (รวมกับเบี้ยประกันชีวิตแล้วไม่เกิน 100,000 บาท) \\
\hline
RMF (ข้อ 2(55)) & ไม่เกินร้อยละ 30 ของเงินได้พึงประเมิน และไม่เกิน 500,000 บาท \\
\hline
SSF (ข้อ 2(102)) & ไม่เกินร้อยละ 30 ของเงินได้พึงประเมิน และไม่เกิน 200,000 บาท \\
\hline
Thai ESG (ข้อ 2(105)) & ไม่เกินร้อยละ 30 ของเงินได้พึงประเมิน และไม่เกิน 100,000--300,000 บาท \\
\hline
Thai ESGX (ข้อ 2(107)) & เงินใหม่ไม่เกิน 300,000 บาท; LTF เดิมที่สับเปลี่ยนไม่เกิน 500,000 บาท \\
\hline
ดอกเบี้ยเงินกู้ยืมเพื่อที่อยู่อาศัย (มาตรา 47(1)(ซ) + ข้อ 2(53)/(52)/(59)) & ไม่เกิน 100,000 บาท (10,000 + 90,000) รวมทุกแหล่งกู้ \\
\hline
เงินสมทบกองทุนประกันสังคม (มาตรา 47(1)(ฌ)) & ตามจำนวนที่จ่ายจริง (สูงสุด 750 บาทต่อเดือนที่อัตราปกติ) \\
\hline
อุปการะเลี้ยงดูบิดามารดา (มาตรา 47(1)(ญ)) & คนละ 30,000 บาท (สูงสุด 4 คน เท่ากับ 120,000 บาท) \\
\hline
เบี้ยประกันสุขภาพบิดามารดา (ข้อ 2(76)) & รวมไม่เกิน 15,000 บาท \\
\hline
อุปการะเลี้ยงดูคนพิการ/คนทุพพลภาพ (มาตรา 47(1)(ฎ)) & คนละ 60,000 บาท \\
\hline
เงินบริจาคพรรคการเมือง (มาตรา 47(1)(ฏ)) & ไม่เกิน 10,000 บาท \\
\hline
เงินบริจาคทั่วไป (มาตรา 47(7)) & ไม่เกินร้อยละ 10 ของเงินได้พึงประเมินหลังหักค่าใช้จ่ายและค่าลดหย่อน \\
\hline
ยกเว้นภาษีผู้สูงอายุ 65 ปีขึ้นไป (ข้อ 2(72)) & ไม่เกิน 190,000 บาท \\
\hline
ยกเว้นภาษีคนพิการอายุไม่เกิน 65 ปี (ข้อ 2(81)) & ไม่เกิน 190,000 บาท (เลือกใช้สิทธิร่วมกับข้อ 2(72) ในปีเดียวกันไม่ได้) \\
\hline
ระบบผลิตไฟฟ้าจากพลังงานแสงอาทิตย์ (พ.ร.ฎ. ฉบับที่ 805) & ไม่เกิน 200,000 บาท \\
\hline
\end{longtable}

\section{บทสรุปและข้อสังเกตท้ายเรื่อง}

จากเนื้อหาทั้งหมดที่กล่าวมา จะเห็นได้ว่าโครงสร้างค่าลดหย่อนภาษีเงินได้บุคคลธรรมดาของประเทศไทยมีความซับซ้อนเชิงระบบในหลายระดับ ดังนี้

\begin{enumerate}
    \item \textbf{ระดับฐานกฎหมาย:} ค่าลดหย่อนหลักส่วนใหญ่มีฐานซ้อนกันสองชั้น คือบทบัญญัติหลักในมาตรา 47 แห่งประมวลรัษฎากร และส่วนเพิ่มเติมในกฎกระทรวงฉบับที่ 126 (เช่น เบี้ยประกันชีวิต เงินสะสมกองทุนสำรองเลี้ยงชีพ และดอกเบี้ยเงินกู้บ้าน)
    \item \textbf{ระดับรูปแบบทางกฎหมาย:} ต้องแยกแยะระหว่าง "ค่าลดหย่อน" ตามมาตรา 47 โดยตรง กับ "ข้อยกเว้นภาษีเงินได้เท่าที่ได้จ่าย" ตามพระราชกฤษฎีกาหรือกฎกระทรวงอื่น ซึ่งมีผลในทางปฏิบัติคล้ายค่าลดหย่อนแต่ที่มาทางกฎหมายต่างกัน
    \item \textbf{ระดับเงื่อนไขเรื่องถิ่นที่อยู่:} ค่าลดหย่อนส่วนตัวไม่มีเงื่อนไขเรื่องถิ่นที่อยู่ แต่ค่าลดหย่อนคู่สมรส บุตร และบิดามารดา มีเงื่อนไขจำกัดตามมาตรา 47(3) ในกรณีผู้มีเงินได้มิได้เป็นผู้อยู่ในประเทศไทย
    \item \textbf{ระดับเงื่อนไขกรณีสามีภริยาต่างมีเงินได้:} ค่าลดหย่อนคู่สมรสมีเพดานรวมจำกัดที่ 120,000 บาทเสมอ (ตามมาตรา 47(2)) แต่ค่าลดหย่อนบุตรไม่มีเงื่อนไขจำกัดเช่นนั้น ทั้งสองฝ่ายหักได้เต็มจำนวนเสมอไม่ว่าจะเลือกวิธียื่นแบบใด
    \item \textbf{ระดับเพดานรวมข้ามประเภท:} ค่าลดหย่อนเพื่อการเกษียณ (กองทุนสำรองเลี้ยงชีพ กบข. กองทุนสงเคราะห์โรงเรียนเอกชน กองทุนการออมแห่งชาติ RMF และประกันชีวิตแบบบำนาญ) มีเพดานรวมกันไม่เกิน 500,000 บาทเสมอ แม้แต่ละประเภทจะมีเพดานเฉพาะของตนเองก็ตาม
\end{enumerate}

ผู้ศึกษาควรจดจำทั้งตัวเลขเพดานของแต่ละรายการ และความสัมพันธ์เชิงเพดานรวมระหว่างรายการต่าง ๆ เพื่อสามารถนำไปประยุกต์ใช้ในการคำนวณภาษีเงินได้บุคคลธรรมดาในสถานการณ์ที่ซับซ้อนได้อย่างถูกต้อง

\end{document}
